% ContEst.tex -- Control and Estimation Co-Design
% Automatica style (autart.cls, two-column). Compile with the autart
% class from the Automatica sample. Per Automatica policy, NO author
% macros are defined in this file.
%
% Bibliography keys come from:
%   (i)  your References.bib:
%        garcia2019control, feldbaum1960dual, bar1974dual,
%        kumar2015stochastic, resnick2019probability, anderson2012optimal
%   (ii) ContEst_related_work.bib (generated alongside this file)
%   (iii) two works you own; suggested BibTeX in the comment block at the
%         end:  grigoriadis1998integrate (paper 2),
%               ramadan_ekfkoopman       (paper 3, eKF--Koopman).
% Make sure every key resolves before running bibtex.

\documentclass[twocolumn,amsthm]{autart}

\usepackage[utf8]{inputenc}
\usepackage{comment}
\usepackage{algorithm}
\newtheorem{definition}{Definition}[section]
\usepackage{algpseudocode}
\newcommand{\Y}{\mathcal{Y}}
\newcommand{\Ll}{\mathcal{L}}
\newcommand{\E}{\mathbb{E\,}}
\newcommand{\Cov}{\textbf{Cov \,}}
\usepackage{graphicx}
\usepackage[version=4]{mhchem}
\usepackage{siunitx}
\usepackage{longtable,tabularx}
\usepackage{booktabs}
\usepackage{pifont}
\setlength\LTleft{0pt} 
\newcommand{\R}{\mathbb{R}}
\usepackage{graphics} % for pdf, bitmapped graphics files
\usepackage{float}
\usepackage{epsfig} % for postscript graphics files
%\usepackage{mathptmx} % assumes new font selection scheme installed
\usepackage{times} % assumes new font selection scheme installed
\usepackage{amsmath} % assumes amsmath package installed
\usepackage{amssymb}  % assumes amsmath package installed
\usepackage{mathtools}% http://ctan.org/pkg/mathtools
\usepackage{color}
\usepackage[dvipsnames]{xcolor}
\usepackage{hyperref}

\newtheorem{assumption}{Assumption}
\newtheorem{lemma}{Lemma}
\newtheorem{theorem}{Theorem}
\newtheorem{proposition}{Proposition}
\newtheorem{corollary}{Corollary}
%\hypersetup{
%    colorlinks=true,
%    linkcolor=blue,
%    filecolor=magenta,      
%    urlcolor=blue,
%}

\usepackage{bm}
\usepackage{times}
% Bob's additions
\usepackage{soul}
\newcommand{\cut}[1]{{\color{blue} \st{#1}}} % replace with {} to cut for real
\newcommand{\add}[1]{{\color{red} #1}} % replace blue with black to add for real
\newcommand{\matlab}{\mbox{\textsc{Matlab}}}

\begin{document}

\begin{frontmatter}

\title{Control and Estimation Co-Design\thanksref{footnoteinfo}}

 \thanks[footnoteinfo]{Corresponding author M.~S.~Ramadan. This material was
 based upon work supported by the U.S.\ Department of Energy, Office of Science,
 Office of Advanced Scientific Computing Research (ASCR), under Contract
 DE-AC02-06CH11347.}

\author[anl]{Mohammad S. Ramadan}\ead{mramadan@anl.gov},
\author[anl]{Philip Dinenis}\ead{pdinenis@anl.gov},
\author[anl]{Mihai Anitescu}\ead{anitescu@mcs.anl.gov}

\address[anl]{Mathematics and Computer Science Division, Argonne National
Laboratory, Lemont, IL 60439, USA}

\begin{keyword}
Control co-design; Stochastic optimal control; State estimation;
Semidefinite programming; Sensor placement; Sensitivity analysis.
\end{keyword}

\begin{abstract}
We extend control co-design to incorporate state and parameter estimation,
under the umbrella of stochastic optimal control (dual control). The proposed
framework, Control and Estimation Co-Design (ContEst), parametrizes the plant
dynamics, the output operator, and the noise covariances to represent design
choices such as actuator tuning, sensor placement, measurement configuration, and
the tuning of the estimator itself; these parameters
affect both closed-loop performance and estimator accuracy. In place of the
sequential plant--control--estimator pipeline, which is in general suboptimal,
ContEst poses a single two-stage problem solved by first-order methods. A key
enabler is that the gradient of each inner control--estimation value with
respect to the design parameters is available exactly from the dual variables
the inner solver already returns, by the envelope theorem; no differentiation
through the optimizer or its optimality conditions is needed. We instantiate the
framework in a linear--quadratic--Gaussian regime, where the control and
estimation costs become two semidefinite programs whose design gradients are
read from the duals of Lyapunov constraints; in a robust ($H_\infty$)
instantiation, where the same two axes become worst-case control and filtering
problems solved by a pair of game Riccati equations sharing the identical
envelope gradient; and in a locally linearized regime, where an extended Kalman
filter yields a convex information-state model predictive control problem whose
design gradients are read from the dynamics costates. We also show how sensor
allocation, a discrete (mixed-integer) selection, admits a convex $\ell_1$
relaxation that promotes sparse sensing.
\end{abstract}

\end{frontmatter}

%=======================================================================
\section{Introduction}\label{sec:intro}
%=======================================================================

Engineering systems are still, overwhelmingly, designed \emph{sequentially}. A
plant is conceived and sized by one set of disciplines; the sensing and
instrumentation are specified next; and the controller and estimator are
designed last, once the plant and the measurement chain are frozen
\cite{garcia2019control}. Each stage constrains the next, so decisions made
early---before any closed-loop consideration---silently bound the performance
achievable at the end. This division of labor is convenient and entrenched, but
it is in general \emph{not optimal}: the sequential solution is one feasible
point of a larger joint problem and need not be a stationary point of it, because
the plant and controller optimizations are \emph{coupled}~\cite{fathy2001coupling},
which motivates solving them jointly~\cite{garcia2019control}.

\emph{Control co-design} (CCD) addresses this on the control side, optimizing the
plant and the controller together~\cite{garcia2019control,allison2014multidisciplinary}---%
either by \emph{nested} strategies (an inner controller solve inside an outer plant
search) or by \emph{simultaneous} ones that optimize both at once, contrasted
in~\cite{herber2019nested}. The same idea appeared in
process systems engineering as \emph{integrated design and control}
\cite{sakizlis2004recent,sharifzadeh2013integration} and in aerospace as
\emph{simultaneous structure/control design}
\cite{onoda1987approach,hale1985optimal}. What these literatures share is a
plant-plus-controller scope: the \emph{sensing} and the \emph{estimator} are
typically fixed, or reduced to a separate sensor-selection step. Yet for any
system that is not fully and noiselessly measured---essentially every real
system---achievable performance is limited by what the estimator can infer,
which is set by where and how well one senses. This echoes the classical
input--output controllability principle that the best attainable closed-loop
performance is fixed by the plant together with the sensor and actuator
configuration, and cannot be recovered by the controller alone
\cite{skogestad2005multivariable}. Sensor placement, precision, and
configuration are design variables on the same footing as actuator authority,
and they couple to the controller through the information available for
feedback.

This paper proposes \emph{Control and Estimation Co-Design} (ContEst): the
design parameters governing actuation \emph{and} sensing/estimation are placed
in one two-stage optimization, whose second-stage objective is the stochastic
optimal control (dual control) value \eqref{eq:Jstoch}---the notion, originating
with Feldbaum, that the optimal input must at once regulate the state and probe to
reduce its uncertainty~\cite{feldbaum1960dual,bar1974dual,mesbah2018stochastic}. Because that cost is
stochastic, it depends on the system not through the unobserved state but through
its \emph{information state}---the filtering density
\cite{kumar2015stochastic}---whose evolution explicitly involves the output
equation. Optimizing it therefore improves the controller, the estimator, and
their interaction at once. A caveat we make precise in
Section~\ref{sec:ekf}: we reserve ``dual control'' for the exact value
\eqref{eq:Jstoch}, whose input has the Bar-Shalom--Tse \emph{dual effect} of
acting on the state \emph{and} on the information gathered
\cite{bar1974dual,mesbah2018stochastic}; the tractable eKF surrogate we optimize
in the nonlinear regime is \emph{caution-aware} (it prices estimation uncertainty
through $\operatorname{tr}(Q\Sigma_{t\mid t})$) but, after linearizing the
information dynamics, \emph{certainty-equivalent in the mean}---it captures
\emph{caution} (hedging against estimation uncertainty) but not \emph{active
probing} (injecting inputs deliberately to reduce it). The very prevalence of sequential
design is thus an opportunity: many existing systems can be revisited and
improved by a joint design pass, often without changing the hardware menu, only
how its continuous parameters are chosen.

\noindent\textbf{Contributions.}~
(i)~We formulate co-design of control \emph{and} estimation as a single two-stage
problem (Problem~\ref{prob:bilevel})---a design first stage over a stochastic
control recourse (the second-stage control response)---expressed through the
information state, so
the sensing operator enters the objective explicitly
(Section~\ref{sec:formulation}). (ii)~We show that the design gradient of each
inner value function is obtained \emph{exactly} from the dual variables the inner
convex solver already returns, by the envelope theorem
(Theorem~\ref{prop:envelope}, Section~\ref{sec:envelope})---which gives the
gradient of an optimal-value function from the optimizer alone, without
differentiating through it; no implicit
differentiation and no differentiation of the Karush--Kuhn--Tucker (KKT) system
is needed---only cheap model Jacobians. (iii)~In the linear--quadratic--Gaussian
(LQG) regime we reduce the inner problem, via the classical separation of
estimation and control, to a control and an
estimation semidefinite program (SDP), with gradients read from the duals of the
Lyapunov constraints (Section~\ref{sec:lqg}). (iv)~In the locally linearized
regime an extended Kalman filter (eKF) lifts the problem to an information state
and a convex model predictive control (MPC) problem, with gradients read from the
dynamics costates (adjoint variables; Section~\ref{sec:ekf}). (v)~We treat sensor allocation as a
mixed-integer selection and show its $\ell_1$ relaxation promotes sparse sensing
(Section~\ref{sec:sensor}). Section~\ref{sec:algorithm} states the outer
algorithm and Section~\ref{sec:numerical} four co-design studies spanning
spacecraft, process, and power systems.

Throughout, $\| x \|_Q^2 = x^\top Q x$, $\mathbb{E}$ is expectation,
$\operatorname{tr}(\cdot)$ the trace, and $M \succeq 0$ ($M \succ 0$) positive
(semi)definiteness. The information available at time $k$ is
$\mathcal{Z}_k = (y_{0:k}, u_{0:k-1})$.

%=======================================================================
\section{Background and related work}\label{sec:related}
%=======================================================================

\emph{Control co-design and integrated design.}
\cite{fathy2001coupling} established that plant and controller optimization are
coupled and cannot be separated without loss; \cite{allison2014multidisciplinary}
cast dynamic-system CCD as multidisciplinary design optimization (MDO); and
\cite{herber2019nested} systematized the nested and simultaneous solution
strategies. In process systems
engineering the same idea was developed largely as mixed-integer dynamic
optimization under uncertainty \cite{diangelakis2017process},
built on Biegler's simultaneous full-discretization approach
\cite{biegler2010nonlinear} and reviewed in
\cite{sakizlis2004recent,sharifzadeh2013integration}. In all of
these the estimator and sensing layer are typically assumed fixed.

\emph{Structural co-design and information architecture.}
The structural-control community produced the linear-matrix-inequality (LMI)
machinery on which our linear regime rests. Following early
simultaneous structure/control designs \cite{onoda1987approach,hale1985optimal},
\cite{grigoriadis1998integrate} alternated a control LMI with a
plant/actuator-parameter LMI---a direct ancestor of our control SDP. Closest to
the present work, \cite{li2008integrating} treated the \emph{information
architecture} (sensor/actuator gains) as decision variables jointly with control
or estimation via LMIs, and \cite{saraf2017h2,das2017sparse} pushed this toward
sparse $H_2$ sensing with guaranteed estimation-error bounds.

\emph{Sensor selection and LQG sensing co-design.}
A large literature treats sensor (and actuator) selection \emph{for estimation}:
convex relaxation \cite{joshi2009sensor}, submodularity and greedy algorithms for
Kalman filtering and network observability
\cite{tzoumas2016sensor,zhang2017sensor}, and
sparsity-promoting joint selection \cite{dhingra2014admm}. These
covariance-based criteria are the control-theoretic face of optimal experimental
design, where information is measured by convex functions of the error covariance
or Fisher information; our estimation cost
$\operatorname{tr}(Q\Sigma^\varepsilon)$ is of A-optimal type. The identification
community's \emph{least-costly experiment design}
\cite{gevers2005identification}---minimizing experiment cost
subject to a closed-loop performance requirement---is the direct sibling of
designing the noise level $V(\theta)$ against closed-loop cost, which our
weight/noise co-design (Section~\ref{sec:weights}) does inside the loop rather than
as a separate identification stage. The closest work to ours
is \cite{tzoumas2021lqg}, who co-design sensing with LQG control through the
separation structure; see also the joint scheduling separation result of
\cite{siami2021separation} and the compositional view of \cite{zardini2021codesign}.
These typically choose \emph{which} sensors to deploy from a discrete menu, in
the linear/LQG regime.

\emph{Dual control.}
That estimation belongs inside the control objective is the content of dual
control \cite{feldbaum1960dual}: an input has a \emph{dual effect}, acting on the
state and on the information gathered about it \cite{bar1974dual}. The
exploration--exploitation reading is surveyed in \cite{mesbah2018stochastic}.

\emph{Differentiable optimization and the envelope theorem.}
Our gradient computation is classical sensitivity analysis
\cite{bonnans2000perturbation} and the envelope theorem \cite{milgrom2002envelope}
specialized to convex programs \cite{boyd2004convex}. Recent work differentiates
through convex programs by implicitly differentiating the KKT/cone solution map
\cite{amos2017optnet,agrawal2019differentiable}, including for bilevel problems
\cite{gould2016differentiating}. As we stress in
Section~\ref{sec:envelope}, for the \emph{value} (not the minimizer) this
machinery is unnecessary---and the difference is quantitative. Implicit
differentiation of the KKT/cone map solves an $m\times m$ linear system per design
direction, with $m=O(r_x^2)$ for the covariance SDPs, i.e.\ $O(r_x^6)$ per
directional derivative, and it differentiates the \emph{minimizer}, which ContEst
never needs; the envelope route reads the already-computed dual and contracts at
$O(r_x^2 r_\theta)$ for the \emph{whole} gradient (Section~\ref{sec:sim:complexity}).

\emph{Positioning: closest work and what is new.}
Three lines of work come closest, and ContEst departs from each in the same
structural way. (i)~\emph{Sensor and estimator design}
\cite{joshi2009sensor,tzoumas2016sensor,zhang2017sensor} places or
selects sensors to sharpen a Kalman-filtering or observability criterion, but for
a \emph{fixed} plant and controller: the actuation and the closed-loop control
are not co-optimized with the sensing. (ii)~\emph{Plant--controller co-design}
\cite{fathy2001coupling,herber2019nested,grigoriadis1998integrate} jointly tunes
plant and controller, but treats sensing and estimation as given---full-state
feedback or a fixed observer---so the output map never enters the design.
(iii)~\emph{Dual and adaptive control}
\cite{feldbaum1960dual,bar1974dual,mesbah2018stochastic} couples
control with online estimation and probing, but for a \emph{fixed} physical and
sensing design $\theta$: the distinction is \emph{offline hardware vs.\ online
policy}---dual/adaptive MPC adapts the \emph{inputs} for a fixed design, whereas
ContEst optimizes the \emph{design} against the (approximated) closed-loop
information dynamics. The nearest single reference is LQG sensing co-design
\cite{tzoumas2021lqg}, which co-selects sensors with an LQG controller through
separation, but only in the linear/LQG regime and as a discrete menu selection;
Section~\ref{sec:sim:pll} compares ContEst head-to-head against the
separation-style sequential pipeline that such menu methods induce.

Table~\ref{tab:positioning} makes these deltas auditable per capability.

\begin{table*}[t]
\centering
\caption{Positioning of ContEst against the nearest lines of work, by capability
(\checkmark\ present, \ding{55}\ absent, \emph{partial} limited). ``Estimator
itself designed'' means tuning a \emph{dynamic} estimator (e.g.\ a PLL bandwidth,
or the noise level $V(\theta)$), not only choosing sensors.}
\label{tab:positioning}
\small
\resizebox{\textwidth}{!}{%
\begin{tabular}{lccccccc}
\toprule
capability & CCD/MDO & sensor & LQG sensing & info-arch. & dual/adaptive
& diff.\ opt. & \textbf{ContEst}\\
 & \cite{garcia2019control,herber2019nested} & selection
 \cite{joshi2009sensor,tzoumas2016sensor} & co-design \cite{tzoumas2021lqg}
 & LMIs \cite{li2008integrating} & MPC \cite{mesbah2018stochastic}
 & layers \cite{amos2017optnet,agrawal2019differentiable} & \\
\midrule
plant / actuation designed & \checkmark & \ding{55} & \ding{55} & partial & \ding{55} & \ding{55} & \checkmark\\
sensing / output map designed & \ding{55} & \checkmark & \checkmark~(menu) & \checkmark & \ding{55} & \ding{55} & \checkmark~(cont.\ + $\ell_1$ menu)\\
estimator itself designed & \ding{55} & \ding{55} & \ding{55} & \ding{55} & \ding{55} & \ding{55} & \checkmark\\
closed-loop control in objective & \checkmark & \ding{55} & \checkmark & partial & \checkmark~(online) & \checkmark & \checkmark\\
nonlinear regime & \checkmark & \ding{55} & \ding{55} & \ding{55} & \checkmark & \checkmark & \checkmark~(eKF surrogate)\\
gradient w.r.t.\ design & often FD & greedy/relax. & menu enum. & alt.\ LMIs & n/a & implicit KKT diff. & exact from duals\\
\bottomrule
\end{tabular}}
\end{table*}

Against these, ContEst is, to our knowledge, the first to close all three loops
at once---plant, controller, \emph{and} estimator---inside a single stochastic
(dual-control) objective expressed through the information state, so that
estimation quality is \emph{optimized} rather than assumed. Its distinguishing
ingredients are: actuation \emph{and} sensing/estimation parameters co-designed as
\emph{continuous} variables in one problem (with discrete sensor allocation
recovered by a sparse $\ell_1$ relaxation, Section~\ref{sec:sensor}); a nonlinear
extension via the eKF and certainty equivalence \cite{ramadan2024extended}; and
design gradients obtained \emph{exactly} from the inner dual variables by the
envelope theorem, without implicit differentiation of the solution map and
without differentiating the KKT system (Section~\ref{sec:envelope}).

\emph{Significance.} Since essentially every deployed system was designed
sequentially, it sits at a feasible but generally non-stationary point of this
joint problem; a single co-design pass can therefore improve it---often without
altering the hardware menu, only how its continuous parameters are chosen. The
studies of Section~\ref{sec:numerical} bear this out, with total-cost reductions
of $6$--$43\%$ in which the sensing axis $\theta_h$---absent from control-only
co-design---is a major, and sometimes dominant, contributor, up to and including
the tuning of the estimator itself (a phase-locked loop, Section~\ref{sec:sim:pll}).

%=======================================================================
\section{Problem formulation}\label{sec:formulation}
%=======================================================================

Consider the parametrized stochastic system
\begin{subequations}\label{eq:ss}
\begin{align}
x_{k+1} &= f(x_k, u_k, w_k;\, \theta), \label{eq:ss-state}\\
y_k     &= h(x_k, v_k;\, \theta), \label{eq:ss-output}
\end{align}
\end{subequations}
with state $x_k \in \mathbb{R}^{r_x}$, input $u_k \in \mathbb{R}^{r_u}$, and
output $y_k \in \mathbb{R}^{r_y}$. The maps $f, h$ are parametrized by
$\theta \in \mathbb{R}^{r_\theta}$, representing plant parameters, actuator
gains, and sensor tuning/placement. The disturbances $w_k, v_k$ are i.i.d.,
mutually independent and independent of $x_0 \sim p_0$ (the prior, before any
measurement including $y_0$).

The goal is to minimize design, control, and estimation costs together. The
\emph{design cost} $J_{\mathrm{des}}(\theta)$ and design set $\theta \in \Theta$
encode the price, manufacturability, or placement difficulty of a $\theta$. The
\emph{control and estimation cost} is the stochastic optimal control value
\begin{equation}\label{eq:Jstoch}
J_{\mathrm{stoch}}(p_0;\theta) = \min_{u_{0:N-1}} \frac{1}{N}\,
\mathbb{E}\Big\{ \sum_{k=0}^{N} \ell(x_k,u_k) + \ell^N(x_N) \Big\},
\end{equation}
subject to constraints $u_k \in \mathbb{U}$, $x_k \in \mathbb{X}$, where
the expectation is over $(x_0, w_{0:N-1}, v_{0:N})$ and $\ell, \ell^N$ are stage
and terminal costs. In general
$x_k \in \mathbb{X}$ is a chance constraint $\mathbf{P}(x_k\in\mathbb{X})\ge
1-\delta$; in this paper we impose the input constraints and the mean-state
constraints in their deterministic (nominal) form within the inner convex
programs of Sections~\ref{sec:lqg}--\ref{sec:ekf} and the studies of
Section~\ref{sec:numerical}, and defer a full probabilistic treatment---which
enters the inner problem as a convex second-order-cone tightening of
$\mathbb{X}$ and leaves the envelope machinery of Section~\ref{sec:envelope}
intact---to future work (Section~\ref{sec:conclusion}).

\begin{prob}[ContEst]\label{prob:bilevel}
\begin{align*}
\textbf{Design (first stage):}\ &\min_{\theta \in \Theta}\
J(\theta) = J_{\mathrm{des}}(\theta) + J_{\mathrm{stoch}}(p_0;\theta),\\
\textbf{Control (second stage):}\ & J_{\mathrm{stoch}}(p_0;\theta)\ \text{as in
\eqref{eq:Jstoch}}.
\end{align*}
\end{prob}

\begin{rem}[Operating scenarios]\label{rem:scenarios}
For clarity Problem~\ref{prob:bilevel} is stated for a single prior $p_0$. Multiple
\emph{operating scenarios}---distinct initial-condition densities $p_0^{(s)}$, load
levels, disturbance statistics, or tracking set-points---are accommodated verbatim
by replacing the second-stage value with a weighted average
\begin{equation}\label{eq:scenario-avg}
J_{\mathrm{stoch}}(\theta)=\sum_{s=1}^{M}\pi_s\,J_{\mathrm{stoch}}(p_0^{(s)};\theta),
\qquad \pi_s\ge0,\ \textstyle\sum_{s=1}^{M}\pi_s=1,
\end{equation}
over the $M$ scenarios with weights $\pi_s$ (e.g.\ $\pi_s=1/M$ for a plain
average, or empirical probabilities of each regime). Each
$J_{\mathrm{stoch}}(p_0^{(s)};\theta)$ is an independent second-stage solve and the
design objective is linear in these values, so the envelope gradient
(Section~\ref{sec:envelope}) is the same convex combination
$\sum_s\pi_s\,\partial_\theta J_{\mathrm{stoch}}(p_0^{(s)};\theta)$ of the
per-scenario gradients: nothing in the machinery changes, only the number of inner
solves per design. The studies of Section~\ref{sec:numerical} use a single scenario
($M=1$).
\end{rem}

\emph{From bilevel to two-stage.} Problem~\ref{prob:bilevel} is, in general form,
a \emph{bilevel} program: an outer optimization over $\theta$ whose objective
contains the argmin of an inner optimization over the control. Solving a bilevel
problem directly is hard---it couples the two levels through the inner
\emph{minimizer} $u^\star(\theta)$, whose sensitivity $\partial u^\star/\partial
\theta$ requires implicitly differentiating the inner optimality conditions (an
MPEC). The structure here, however, collapses this: the design $\theta$ is a
\emph{here-and-now} decision fixed before any realization, and the control is the
\emph{recourse} that reacts to the information $\mathcal{Z}_k$, so
Problem~\ref{prob:bilevel} is precisely a \emph{two-stage stochastic program}
with first-stage variable $\theta$ and second-stage (recourse) value
$J_{\mathrm{stoch}}(p_0;\theta)$. What the outer method needs from the second
stage is not its minimizer but only the sensitivity of its \emph{value}, and for
that the envelope theorem (Section~\ref{sec:envelope}) supplies the gradient
directly from the duals the inner solver already returns---no differentiation of
$u^\star(\theta)$ and no MPEC. It is this reduction from a bilevel to a two-stage
problem that makes ContEst tractable, and we treat it as two-stage throughout.

Compared with finite-difference sensitivities, ContEst obtains the full design
gradient at essentially the cost of a single inner solve: forward (central)
differences need $r_\theta+1$ ($2r_\theta$) inner solves for an $r_\theta$-vector
$\theta$, whereas ContEst solves the inner problem once and contracts its returned
duals or costates against model Jacobians---an $O(r_\theta)$ speedup---and returns
the \emph{exact} gradient of the inner value rather than a truncated, step-size
sensitive approximation.

\subsection{The information state}
The second stage is a stochastic (dual) control problem
\cite{feldbaum1960dual,bar1974dual}. A causal law must be a function of the
available information $\mathcal{Z}_k$, not of the unobserved $x_k$; the
sufficient statistic is the filtering density $p(x_k \mid \mathcal{Z}_k;\theta)$,
governed by the Bayesian filter. With $\mathcal{Z}_k = (y_{0:k}, u_{0:k-1})$, the
filter alternates a time update,
\begin{equation}\label{eq:Tupdate}
p(x_{k+1} \mid \mathcal{Z}_k, u_k) = \int p(x_{k+1} \mid x_k, u_k;\theta)\,
p(x_k \mid \mathcal{Z}_k)\, dx_k,
\end{equation}
and a measurement update,
\begin{equation}\label{eq:Mupdate}
p(x_{k+1} \mid \mathcal{Z}_{k+1}) \propto p(y_{k+1} \mid x_{k+1};\theta)\,
p(x_{k+1} \mid \mathcal{Z}_k, u_k).
\end{equation}
Both \eqref{eq:Tupdate}--\eqref{eq:Mupdate} involve the output
equation~\eqref{eq:ss-output} and hence the sensing parameters in $\theta$. The
standard smoothing identity makes this explicit in the cost.

\begin{lem}[Smoothing {\cite[Ch.~10]{resnick2019probability}}]\label{lem:smoothing}
For a probability space $(\Omega,\mathcal{A},\mathbf{P})$, a measurable
$\chi:\Omega\to\mathbb{R}$ with $\mathbb{E}_{\mathbf P}|\chi|<\infty$, and
$\sigma$-algebras $\mathcal{A}_0 \subset \mathcal{A}_1 \subset \mathcal{A}$,
$\mathbb{E}_{\mathbf P}\{\chi \mid \mathcal{A}_1\} =
\mathbb{E}_{\mathbf P}\{\mathbb{E}_{\mathbf P}\{\chi \mid \mathcal{A}_0\} \mid
\mathcal{A}_1\}$.
\end{lem}

\begin{lem}\label{lem:stagecost}
For a causal law $u_k = u_k(\mathcal{Z}_k, p_0)$,
\begin{equation*}
\mathbb{E}\,\ell(x_k,u_k) = \!\iint \ell(x_k,u_k)\,
p(x_k \mid \mathcal{Z}_k;\theta)\, p(\mathcal{Z}_k;\theta)\, dx_k\, d\mathcal{Z}_k.
\end{equation*}
\end{lem}
\begin{proof}
Conditioning on $\mathcal{Z}_k$ yields a $\sigma$-algebra contained in the one
underlying \eqref{eq:Jstoch}; apply Lemma~\ref{lem:smoothing}.
\end{proof}

Lemma~\ref{lem:stagecost} shows the cost depends on the anticipated filtering
densities, hence on the output equation: optimizing it improves the controller,
the filter, and their coupling. (Causality requires conditioning on
$\mathcal{Z}_{k'}$ with $k' \le k$; we take $k' = k$.)

%=======================================================================
\section{Exact design gradients via the envelope theorem}\label{sec:envelope}
%=======================================================================

The two-stage problem is solved by descent on the first-stage variable $\theta$,
so the crux is $\nabla_\theta J_{\mathrm{stoch}}$: the sensitivity of a
\emph{value function} to parameters entering its objective and constraints. Our
central observation is that this gradient is essentially free---and it is exactly
this that reduces the bilevel program to the two-stage one. Write the convex
second-stage problem in conic form,
\begin{equation}\label{eq:generic}
V(\theta) = \min_{z}\ \phi(z,\theta)\quad \text{s.t.}\quad g(z,\theta)=0,\ \
G(z,\theta) \in \mathcal{K},
\end{equation}
with closed convex cone $\mathcal{K}$ (the positive semidefinite cone for SDPs),
Lagrangian $L(z,\mu,S,\theta) = \phi(z,\theta) + \langle \mu, g(z,\theta)\rangle
+ \langle S, G(z,\theta)\rangle$, and optimal primal--dual triple
$(z^\star,\mu^\star,S^\star)$. We collect the hypotheses under which the value
$V$ is differentiable and its gradient is read from that triple.

\begin{assumption}\label{as:env}
On a neighborhood $\Theta_0$ of the design point $\theta$:
\emph{(A1)}~for each $\theta\in\Theta_0$, \eqref{eq:generic} is convex in $z$,
its data $(\phi,g,G)$ are jointly $C^1$ in $(z,\theta)$, and Slater's condition
holds, so strong duality holds with an attained primal--dual solution;
\emph{(A2)}~the primal minimizer $z^\star(\theta)$ is unique;
\emph{(A3)}~the primal--dual solution is \emph{nondegenerate} and strictly
complementary, so the dual $(\mu^\star,S^\star)$ is unique
\cite{alizadeh1997complementarity,shapiro1997differentiability};
\emph{(A4)}~$\theta$ enters \eqref{eq:generic} only through finitely many data
objects $D(\theta)$ (e.g.\ $A_\theta,B_\theta,C_\theta,W,V$) that are $C^1$
in $\theta$.
\end{assumption}

Assumptions (A1)--(A3) are exactly the regularity under which the optimal value
of a conic (in particular semidefinite) program is differentiable and its
sensitivity is governed by the unique dual
\cite{bonnans2000perturbation,shapiro1997differentiability}; (A4) is a modeling
convenience that isolates the $\theta$-dependence in cheap model Jacobians.

\begin{theorem}[Exact design gradient from the inner dual]\label{prop:envelope}
Under Assumption~\ref{as:env}, $V$ is differentiable at $\theta$ and
\begin{equation}\label{eq:envelope}
\nabla_\theta V(\theta) = \nabla_\theta L(z^\star,\mu^\star,S^\star,\theta)
= \partial_\theta \phi + \langle \mu^\star, \partial_\theta g\rangle
+ \langle S^\star, \partial_\theta G\rangle,
\end{equation}
all partials evaluated at $(z^\star,\theta)$ with $z^\star,\mu^\star,S^\star$
held fixed.
\end{theorem}
\begin{proof}
By (A1)--(A3) the solution map $\theta\mapsto(z^\star,\mu^\star,S^\star)$ is
single-valued and differentiable on $\Theta_0$
\cite{bonnans2000perturbation,shapiro1997differentiability}, and strong duality
with complementary slackness gives
$V(\theta) = L(z^\star(\theta),\mu^\star(\theta),S^\star(\theta),\theta)$.
Differentiating and using inner stationarity $\nabla_z L = 0$ and the
stationarity of $L$ in $(\mu,S)$ at the solution,
$\nabla_\theta V = \nabla_z L\,\partial_\theta z^\star
+ \nabla_{(\mu,S)}L\,\partial_\theta(\mu^\star,S^\star) + \partial_\theta L
= \partial_\theta L$ \cite{milgrom2002envelope,bonnans2000perturbation}.
\end{proof}

\begin{rem}[Degenerate case]\label{rem:degenerate}
If (A3) fails, the dual solution set $\mathcal{D}^\star(\theta)$ is not a
singleton and $V$ is in general only directionally differentiable, with
$V'(\theta;d) = \min_{(\mu,S)\in\mathcal{D}^\star}
\langle \nabla_\theta L(z^\star,\mu,S,\theta), d\rangle$
\cite{bonnans2000perturbation}. The right-hand side of \eqref{eq:envelope}
evaluated at \emph{any} returned dual is then a subgradient rather than the
gradient, and the outer method of Section~\ref{sec:algorithm} degrades
gracefully to a subgradient step. In our studies strict complementarity held at
every solution (consistent with the gradient check of
Section~\ref{sec:numerical}). By Proposition~\ref{prop:lqgreg}, in the
LQG regime degeneracy can arise only where that proposition's hypotheses fail
(e.g.\ a singular process noise $W$, or a loss of detectability at isolated
$\theta$), which pinpoints in advance the design points that need the subgradient
fallback.
\end{rem}

Three points give ContEst its practical edge. \emph{(1)~No implicit
differentiation:} \eqref{eq:envelope} has no term $\partial_\theta z^\star$; we
never compute how the minimizer moves with $\theta$, which is precisely what
implicit-function-theorem methods spend effort on
\cite{amos2017optnet,agrawal2019differentiable}. \emph{(2)~No KKT
differentiation:} the multipliers $(\mu^\star,S^\star)$ are not recomputed by
differentiating optimality conditions; they are the dual solution the inner
solver already returns, and we read them off. \emph{(3)~Only model Jacobians
remain:} in our setting $\theta$ enters \eqref{eq:generic} through a small set of
data objects $D$ (e.g.\ the matrices $A_\theta, B_\theta, \dots$), so
$\nabla_\theta V = \sum_D \langle \partial V / \partial D,\ \partial D /
\partial\theta\rangle$, where $\partial V/\partial D$ comes from
\eqref{eq:envelope} and $\partial D/\partial\theta$ are explicit (or
automatically differentiated) Jacobians of the \emph{model}, not the optimizer.

%=======================================================================
\section{The LQG case: LQR and Kalman SDPs}\label{sec:lqg}
%=======================================================================

Let \eqref{eq:ss} be linear (or linearized about an operating point),
$x_{k+1} = A_\theta x_k + B_\theta u_k + w_k$, $y_k = C_\theta x_k + v_k$, with
zero-mean white $w_k, v_k$ of covariances $W$ and $V$ (possibly functions of
$\theta$). With the observer
$\hat x_{k+1} = A_\theta \hat x_k + B_\theta u_k + G(y_k - C_\theta \hat x_k)$,
$A_\theta - G C_\theta$ stable, the error
$\varepsilon_k = x_k - \hat x_k$ obeys
$\varepsilon_{k+1} = (A_\theta - G C_\theta)\varepsilon_k + w_k - G v_k$.
For $J_\theta = \lim_{k\to\infty} \mathbb{E}\{ x_k^\top Q x_k + u_k^\top R u_k\}$
with $u_k = K\hat x_k$ and $x_k = \hat x_k + \varepsilon_k$, the cyclic property
of the trace gives, with $\Sigma = \mathbf{Cov}(\hat x_k)$ and
$\Sigma^\varepsilon = \mathbf{Cov}(\varepsilon_k)$,
\begin{equation}
J_\theta = \underbrace{\operatorname{tr}(Q\Sigma^\varepsilon)}_{J_{\mathrm{est}}(\theta)}
+ \underbrace{\operatorname{tr}(Q\Sigma) + \operatorname{tr}(R K \Sigma K^\top)}_{J_{\mathrm{cont}}(\theta)}.
\end{equation}
This is the separation principle, and we use it \emph{only within the inner
solve}: for a \emph{fixed} $\theta$, $J_{\mathrm{est}}$ and $J_{\mathrm{cont}}$
depend on disjoint variables and may be minimized independently
\cite{tzoumas2021lqg}. It is \emph{not} a design-level separation. At the design
level $\theta$ couples $A_\theta, B_\theta, C_\theta$ and the noise covariances,
hence both the control gain $K$ and the observer gain $G$, so
$J_{\mathrm{est}}(\theta)$ and $J_{\mathrm{cont}}(\theta)$ vary over a common
$\theta$ and the co-design objective is \emph{not} separable across the two
axes---which is precisely why the joint two-stage formulation of
Problem~\ref{prob:bilevel} is needed rather than a place-then-control pipeline.

\emph{Control SDP.} Following the output-variance-constrained construction
\cite{grigoriadis1998integrate,boyd2004convex},
\begin{equation}\label{eq:Jcont}
\begin{aligned}
J_{\mathrm{cont}}(\theta) = &\min_{\Sigma,L,Z_0}\
\operatorname{tr}(Q\Sigma) + \operatorname{tr}(R Z_0)\\
\text{s.t. } & \Sigma \succ 0, \quad
\left[\begin{smallmatrix} Z_0 & L \\ L^\top & \Sigma \end{smallmatrix}\right]
\succeq 0,\\
& \left[\begin{smallmatrix}
\Sigma - W & A_\theta\Sigma + B_\theta L \\
(A_\theta\Sigma + B_\theta L)^\top & \Sigma
\end{smallmatrix}\right] \succeq 0,
\end{aligned}
\end{equation}
with optimal gain $K^\star = L^\star (\Sigma^\star)^{-1}$.

\emph{Estimation SDP.} The error covariance solves a Lyapunov equation that,
after $F = P^\varepsilon G$ and the LMI relaxation, gives
\begin{equation}\label{eq:Jest}
\begin{aligned}
J_{\mathrm{est}}(\theta) = &\min_{P^\varepsilon,F,Z_1}\
\operatorname{tr}(W P^\varepsilon) + \operatorname{tr}(V Z_1)\\
\text{s.t. } & P^\varepsilon \succ 0, \quad
\left[\begin{smallmatrix} Z_1 & F^\top \\ F & P^\varepsilon \end{smallmatrix}
\right] \succeq 0,\\
& \left[\begin{smallmatrix}
P^\varepsilon - Q & A_\theta^\top P^\varepsilon - C_\theta^\top F^\top \\
P^\varepsilon A_\theta - F C_\theta & P^\varepsilon
\end{smallmatrix}\right] \succeq 0,
\end{aligned}
\end{equation}
dual in form to \eqref{eq:Jcont} and exposing the sensing operator $C_\theta$,
i.e.\ the estimation half of the co-design \cite{li2008integrating}.

\emph{Design gradients from the LMI duals.} Both \eqref{eq:Jcont} and
\eqref{eq:Jest} are instances of \eqref{eq:generic}. Applying
Theorem~\ref{prop:envelope} to the Lyapunov constraint of \eqref{eq:Jcont},
whose dual we partition as
$S = \left[\begin{smallmatrix} S_{11} & S_{12} \\ S_{12}^\top & S_{22}
\end{smallmatrix}\right] \succeq 0$, gives
\begin{equation}\label{eq:gradAB}
\frac{\partial J_{\mathrm{cont}}}{\partial A_\theta} = -2\, S_{12}\Sigma^\star,
\qquad
\frac{\partial J_{\mathrm{cont}}}{\partial B_\theta} = -2\, S_{12}(L^\star)^\top,
\end{equation}
with analogous expressions for $W, Q, R$ from the remaining Lagrangian terms, and
a parallel set for \eqref{eq:Jest}. The design gradient follows by the chain rule
of Theorem~\ref{prop:envelope}(3),
\begin{equation}\label{eq:chain}
\frac{\partial J_{\mathrm{cont}}}{\partial \theta_i} =
\operatorname{tr}\!\Big(\frac{\partial J_{\mathrm{cont}}}{\partial A_\theta}^{\!\top}
\frac{\partial A_\theta}{\partial \theta_i}\Big)
+ \operatorname{tr}\!\Big(\frac{\partial J_{\mathrm{cont}}}{\partial B_\theta}^{\!\top}
\frac{\partial B_\theta}{\partial \theta_i}\Big) + \cdots,
\end{equation}
where the model Jacobians $\partial A_\theta/\partial\theta_i$, etc., are cheap
and exact. No sweep through the SDP solver is required.

Rather than \emph{assume} the regularity of Assumption~\ref{as:env} for these
SDPs, we can \emph{verify} it from the model data. The following makes the
hypotheses checkable and, in doing so, identifies the LMI duals as classical
Riccati objects.

\begin{proposition}[LQG regularity]\label{prop:lqgreg}
Fix $\theta$ and suppose $(A_\theta,B_\theta)$ is stabilizable, $(A_\theta,
C_\theta)$ is detectable, and $Q,R,W,V\succ0$, with all data $C^1$ in $\theta$.
Then for the control SDP~\eqref{eq:Jcont} and the estimation
SDP~\eqref{eq:Jest}: \emph{(a)}~Slater's condition holds; \emph{(b)}~the primal
minimizer is unique---$\Sigma^\star$ the stabilizing closed-loop covariance and
$K^\star=L^\star(\Sigma^\star)^{-1}$ the Riccati gain, and dually
$P^{\varepsilon\star}=\Sigma^\varepsilon$ the Kalman error covariance;
\emph{(c)}~the Lyapunov-LMI dual is unique and strictly complementary---in fact
the $(1,1)$ dual block of the control LMI \textbf{is} the control Riccati
solution $P$ (so $S_{11}^\star=P$) and the $(1,1)$ dual block of the estimation
LMI \textbf{is} the steady-state Kalman error covariance $\Sigma^\varepsilon$ (so
$S_{11}^{\prime\star}=\Sigma^\varepsilon$). Hence Assumption~\ref{as:env} holds
on a neighborhood $\Theta_0$ of $\theta$, and Corollary~\ref{cor:lqg} applies.
\end{proposition}

\emph{Proof sketch.} A strictly feasible point is built from any stabilizing gain
inflated by $\varepsilon I$ (Slater); with $W\succ0$ a Lyapunov-eigenvector
argument forces every feasible $\Sigma$ to encode a stabilizing gain, and the
completion-of-squares identity $J(K)-\operatorname{tr}(PW)=\operatorname{tr}\!\big(
(K-K^\star)^\top(R+B^\top PB)(K-K^\star)\Sigma_K\big)$ gives primal uniqueness;
complementary slackness reduces the dual stationarity to the closed-loop Stein
equation $\Lambda=Q+K^{\star\top}RK^\star+A_{\mathrm{cl}}^\top\Lambda
A_{\mathrm{cl}}$, whose unique solution is $\Lambda=P$, and strict
complementarity follows from rank counting with $P\succeq Q\succ0$, $R\succ0$.
The full argument and the dual identities are in Appendix~\ref{app:lqgreg}.
\hfill$\square$

Controllability and observability imply the stabilizability/detectability
hypotheses, so the practical reading of Proposition~\ref{prop:lqgreg} is:
\emph{controllable, observable plant with positive-definite weights $\Rightarrow$
everything in Sections~\ref{sec:envelope}--\ref{sec:lqg} is exact and
differentiable}. It also explains structurally why the Riccati gradient formulas
of Section~\ref{sec:sim:complexity} coincide with the LMI-dual
formulas~\eqref{eq:gradQRcont}--\eqref{eq:gradQRest}: they are literally the same
object, $S_{11}^\star=P$.

\begin{corollary}[LQG design gradient]\label{cor:lqg}
Under the hypotheses of Proposition~\ref{prop:lqgreg} (stabilizability,
detectability, $Q,R,W,V\succ0$), the SDPs~\eqref{eq:Jcont}--\eqref{eq:Jest}
satisfy Assumption~\ref{as:env}, so the Lyapunov-LMI dual $S$ (resp.\ its
estimation counterpart) is unique. Then $J_{\mathrm{cont}}(\theta)$ and
$J_{\mathrm{est}}(\theta)$ are differentiable at $\theta$ and their design
gradients are given exactly by \eqref{eq:gradAB}--\eqref{eq:chain} (resp.\ the
parallel expressions in $C_\theta, W, V$), read from the returned LMI duals.
\end{corollary}

\subsection{Co-designing the weights and noise levels}\label{sec:weights}

Nothing in the framework restricts $\theta$ to the dynamics and output maps. The
weights $Q, R$ of the quadratic cost and the noise covariances $W, V$ of the
process and measurement channels are simply more of the data objects $D(\theta)$
of Assumption~\ref{as:env}(A4), so making any of them a design variable requires
\emph{no new machinery}: the same duals contract against a new model Jacobian.
This matters in practice because the noise levels are exactly the knobs a sensor
or actuator budget buys---a better inertial unit lowers $V$, a stiffer mount or a
cleaner drive lowers $W$---so ``how good a sensor to buy'' is a continuous
co-design variable on the same footing as ``where to place it.''

Concretely, both LQG subproblems are affine in these data, so their sensitivities
are read off directly, with no LMI-dual partition required for the objective
terms. From the control SDP~\eqref{eq:Jcont},
\begin{equation}\label{eq:gradQRcont}
\frac{\partial J_{\mathrm{cont}}}{\partial Q} = \Sigma^\star,\qquad
\frac{\partial J_{\mathrm{cont}}}{\partial R} = Z_0^\star,\qquad
\frac{\partial J_{\mathrm{cont}}}{\partial W} = -\,S_{11}^\star,
\end{equation}
the first two from the objective $\operatorname{tr}(Q\Sigma)+\operatorname{tr}(RZ_0)$ and the
last from the $(1,1)$ block $\Sigma-W$ of the Lyapunov LMI, whose dual is
$S_{11}^\star$. Dually, from the estimation SDP~\eqref{eq:Jest},
\begin{equation}\label{eq:gradQRest}
\frac{\partial J_{\mathrm{est}}}{\partial W} = P^{\varepsilon\star},\qquad
\frac{\partial J_{\mathrm{est}}}{\partial V} = Z_1^\star,\qquad
\frac{\partial J_{\mathrm{est}}}{\partial Q} = -\,S_{11}^{\prime\star},
\end{equation}
with $S_{11}^{\prime\star}$ the $(1,1)$ dual block of the estimation LMI. Two
structural facts are worth noting. First, $Q$ and $W$ each appear in \emph{both}
subproblems---$Q$ weighs the control cost $\operatorname{tr}(Q\Sigma)$ and the estimation
cost $\operatorname{tr}(Q\Sigma^\varepsilon)$, and $W$ is both the process noise seen by the
estimator and a term in the control Lyapunov constraint---so the total gradient is
just the sum of the two contributions, e.g.\ $\partial J_\theta/\partial W =
-S_{11}^\star + P^{\varepsilon\star}$. Second, these expressions substitute into the
same chain rule~\eqref{eq:chain}: whatever the underlying design variable
$\theta_i$ that moves $Q,R,W$ or $V$, its gradient is
$\operatorname{tr}((\partial J/\partial Q)^\top \partial Q/\partial\theta_i) + \cdots$,
with cheap model Jacobians. The eKF case of Section~\ref{sec:ekf} is identical in
spirit: $Q,R$ enter the stage cost $\ell$ directly, while the noise covariances
$\Sigma_w, \Sigma_v$ enter the filter map and hence
$A_\theta^{\mathrm{info}}, B_\theta^{\mathrm{info}}$, so their gradients are
delivered by the costate contraction~\eqref{eq:gradmpc} with no change.

A caveat keeps the extension honest. Designing $W, V$ is unambiguous---they are
physical noise levels priced by $J_{\mathrm{des}}$, and lowering them always helps
the inner value, so the design cost is what makes the trade-off nontrivial.
Designing $Q, R$ is meaningful only when the weights are genuinely negotiable
knobs---an actuation-effort budget, or a specification the designer is free to
retune---rather than a fixed performance specification; absent a competing term in
$J_{\mathrm{des}}$, minimizing the closed-loop cost over its own weights is
degenerate (e.g.\ $R\to0$). We therefore treat $Q,R$ as design variables only
against an explicit $J_{\mathrm{des}}$ penalty or bound, exactly as for the other
$\theta$ coordinates.

\subsection{The robust analog: $H_\infty$ control and filtering SDPs}
\label{sec:hinf}

The LQG development above prices \emph{average} performance against
\emph{stochastic} white disturbances. When the disturbances are better modeled as
unknown-but-bounded $\ell_2$ signals and the designer cares about the
\emph{worst case}, the same linear plant admits an $H_\infty$ inner problem, and
--- crucially for ContEst --- it again reduces to a pair of semidefinite programs
of the generic form~\eqref{eq:generic}, so Theorem~\ref{prop:envelope} delivers
the design gradient with no new machinery. The inner \emph{value} is now the
squared worst-case closed-loop $\ell_2$-gain $\gamma^2(\theta)$ rather than a
stationary variance, but it is still the optimal value of a conic program whose
$\theta$-dependence is confined to the model data.

Write the linear plant with an explicit disturbance channel and a performance
output,
\begin{equation}\label{eq:hinf_plant}
x_{k+1} = A_\theta x_k + B_\theta u_k + E_\theta w_k,\qquad
z_k = C_{z}\,x_k + D_{z}\,u_k,
\end{equation}
where $w_k\in\ell_2$ is the exogenous disturbance, $E_\theta$ its input map
(e.g.\ $E_\theta = W_\theta^{1/2}$, so the process ``noise level'' of the LQG model
becomes a disturbance-channel weight), and $z_k$ stacks the penalized
state and input, $C_z = [\,Q^{1/2};\,0\,]$, $D_z = [\,0;\,R^{1/2}\,]$, so the
$H_\infty$ objective weighs exactly the same signals as $J_{\mathrm{cont}}$.

\emph{Control SDP.} By the discrete-time bounded-real lemma with the standard
linearizing change of variables $Y = P^{-1}\succ 0$, $L = KY$
\cite{boyd1994linear,grigoriadis1998integrate}, the state-feedback $H_\infty$
synthesis problem is the SDP (with $\mu := \gamma^2$)
\begin{equation}\label{eq:Hinf_cont}
\begin{aligned}
\gamma^2_{\mathrm{cont}}(\theta) = &\min_{Y,L,\mu}\ \mu \\
\text{s.t. } & Y \succ 0,\ \
\begin{bmatrix}
-Y & \star & \star & \star \\
0 & -\mu I & \star & \star \\
A_\theta Y + B_\theta L & E_\theta & -Y & \star \\
C_z Y + D_z L & 0 & 0 & -I
\end{bmatrix} \preceq 0,
\end{aligned}
\end{equation}
where $\star$ denotes the block inferred by symmetry, with optimal gain
$K^\star = L^\star (Y^\star)^{-1}$. This is the exact robust
counterpart of the output-variance-constrained control SDP~\eqref{eq:Jcont}: the
Lyapunov certificate $Y$ plays the role of the covariance $\Sigma$, the extra
disturbance row/column carries the gain bound $\mu$, and the performance block
$C_zY+D_zL$ replaces the traced objective.

\emph{Estimation SDP.} Dually, for the filter
$\hat x_{k+1} = A_\theta \hat x_k + G(y_k - C_\theta \hat x_k)$ with measurement
$y_k = C_\theta x_k + D_{v}\,v_k$ and estimation-error output
$\zeta_k = H(x_k - \hat x_k)$, the bounded-real lemma applied to the error
dynamics --- now affine in $(P,F)$ under $F = P G$ --- gives the $H_\infty$
filtering SDP
\begin{equation}\label{eq:Hinf_est}
\begin{aligned}
\gamma^2_{\mathrm{est}}(\theta) = &\min_{P,F,\mu}\ \mu \\
\text{s.t. } & P \succ 0,\ \
\begin{bmatrix}
-P & \star & \star & \star \\
A_\theta^\top P - C_\theta^\top F^\top & -P & \star & \star \\
E_\theta^\top P - D_v^\top F^\top & 0 & -\mu I & \star \\
0 & H & 0 & -I
\end{bmatrix} \preceq 0,
\end{aligned}
\end{equation}
with optimal gain $G^\star = (P^\star)^{-1} F^\star$, dual in form
to~\eqref{eq:Hinf_cont} and exposing the sensing operator $C_\theta$ and
measurement weight $D_v = V_\theta^{1/2}$ --- the estimation half of the robust
co-design \cite{li2008integrating}. As in the LQG case the fixed-$\theta$
separation is an inner convenience: at the design level $\theta$ couples
$A_\theta, B_\theta, C_\theta, E_\theta, D_v$ across \emph{both} LMIs, so the
worst-case co-design objective is again non-separable and Problem~\ref{prob:bilevel}
is needed.

\emph{Design gradients from the BRL duals.} Both \eqref{eq:Hinf_cont} and
\eqref{eq:Hinf_est} are instances of \eqref{eq:generic}: the constraint is a
single conic inclusion $-\mathcal{M}(\cdot,\theta)\in\mathcal{K}$ (the PSD cone),
so Theorem~\ref{prop:envelope} gives
$\nabla_\theta \gamma^2 = -\langle S^\star, \partial_\theta \mathcal{M}\rangle$,
with $S^\star\succeq 0$ the returned BRL dual and $\partial_\theta\mathcal{M}$
nonzero only in the blocks that carry $\theta$. Partition the control dual
conformally with the $4\times4$ LMI of \eqref{eq:Hinf_cont},
$S^\star = [\,S_{ij}\,]$, $S_{ij}=S_{ji}^\top$. The plant blocks
$A_\theta Y + B_\theta L$ sit in position $(3,1)$, giving
\begin{equation}\label{eq:gradHinf}
\frac{\partial \gamma^2_{\mathrm{cont}}}{\partial A_\theta} = -2\,S_{13}\,Y^\star,
\qquad
\frac{\partial \gamma^2_{\mathrm{cont}}}{\partial B_\theta} = -2\,S_{13}\,(L^\star)^\top,
\end{equation}
with analogous expressions for $C_z, D_z$ (from the $(4,1)$ block $S_{14}$) and for
$E_\theta$ (from the $(3,2)$ block $S_{23}$), and structurally identical to the
LQG control gradient~\eqref{eq:gradAB} with the covariance dual $S_{12}$ replaced
by the BRL dual $S_{13}$. For the filter, partitioning the dual of
\eqref{eq:Hinf_est} as $\tilde S = [\,\tilde S_{ij}\,]$, the off-diagonal block
carrying $PA_\theta - FC_\theta$ yields
\begin{equation}\label{eq:gradHinf_est}
\frac{\partial \gamma^2_{\mathrm{est}}}{\partial A_\theta} = -2\,P^\star \tilde S_{12},
\qquad
\frac{\partial \gamma^2_{\mathrm{est}}}{\partial C_\theta} = 2\,(F^\star)^\top \tilde S_{12},
\end{equation}
with parallel terms in $E_\theta, D_v$. In both cases the design gradient closes
through the same chain rule~\eqref{eq:chain} against cheap model Jacobians
$\partial A_\theta/\partial\theta_i,\ \partial C_\theta/\partial\theta_i,\dots$;
as before, no sweep through the SDP solver is required.

\begin{corollary}[$H_\infty$ design gradient]\label{cor:hinf}
Let \eqref{eq:Hinf_cont} (resp.\ \eqref{eq:Hinf_est}) satisfy
Assumption~\ref{as:env} --- for these SDPs, strict feasibility (a stabilizing gain
achieving a gain strictly below the optimum, i.e.\ Slater) and nondegeneracy of
the optimal solution \cite{alizadeh1997complementarity}, so the BRL dual $S^\star$
(resp.\ $\tilde S^\star$) is unique. Then $\gamma^2_{\mathrm{cont}}(\theta)$ and
$\gamma^2_{\mathrm{est}}(\theta)$ are differentiable at $\theta$ and their design
gradients are given exactly by \eqref{eq:gradHinf} (resp.\ \eqref{eq:gradHinf_est})
and their analogs, read from the returned BRL duals. If strict complementarity
fails, Remark~\ref{rem:degenerate} applies verbatim and the returned dual yields a
subgradient.
\end{corollary}

\begin{rem}[Where the SDP route is exact, and where it is not]\label{rem:hinfreg}
The two regularity regimes must be kept separate. At a \emph{fixed} $\gamma$ above
the optimal attenuation the story is clean and is stated as
Proposition~\ref{prop:hinfreg} below (the fixed-$\gamma$ game Riccati is what the
robust study of Section~\ref{sec:sim:mtdc} actually uses). At the \emph{minimal}
$\gamma=\gamma^\star$ that \eqref{eq:Hinf_cont} minimizes over, by contrast, the
BRL LMI loses rank and strict complementarity \emph{generically fails}, so only
the subgradient statement of Remark~\ref{rem:degenerate} is available; this is why
we never differentiate at $\gamma^\star$ and prefer the fixed-$\gamma$ Riccati
form for numerics.
\end{rem}

\emph{Scalable game-Riccati route.} Exactly as the $H_2$ control SDP~\eqref{eq:Jcont}
has an $O(r_x^3)$ Riccati twin ($J_{\mathrm{cont}}=\operatorname{tr}(PW)$ from the
control DARE, used for the large-scale study of
Section~\ref{sec:sim:mtdc}), the $H_\infty$ control SDP admits a Riccati twin that
avoids the interior-point solve entirely --- important because, at the
\emph{minimal}-$\gamma$ optimum, the BRL LMI \eqref{eq:Hinf_cont} loses rank and its
dual is numerically delicate. Fix the attenuation level $\gamma$ and pose the
soft-constrained dynamic game
$\min_u\max_w\ \sum_k\big(x_k^\top Qx_k+u_k^\top Ru_k-\gamma^2 w_k^\top w_k\big)$
for the plant \eqref{eq:hinf_plant} with $E_\theta=W_\theta^{1/2}$; its value is
$x_0^\top X x_0$ with $X\succeq0$ the stabilizing solution of the discrete
$H_\infty$ game Riccati equation. Writing $\tilde B_\theta=[\,B_\theta\ \ E_\theta\,]$
and the \emph{indefinite} game weight
$\tilde R=\operatorname{diag}(R,-\gamma^2 I)$, that GARE is precisely a discrete
algebraic Riccati equation in the augmented data $(A_\theta,\tilde B_\theta,Q,\tilde R)$,
\begin{equation}\label{eq:gare}
X = Q + A_\theta^\top X A_\theta
 - A_\theta^\top X\tilde B_\theta(\tilde R+\tilde B_\theta^\top X\tilde B_\theta)^{-1}
   \tilde B_\theta^\top X A_\theta,
\end{equation}
admissible when $R+B_\theta^\top X B_\theta\succ0$ and the game closed loop
$A_{\mathrm{cl}}=A_\theta+\tilde B_\theta\tilde K$,
$\tilde K=-(\tilde R+\tilde B_\theta^\top X\tilde B_\theta)^{-1}\tilde B_\theta^\top X A_\theta
=[\,K;\,K_w\,]$, is Schur. The \emph{worst-case (guaranteed) stationary cost}
\begin{equation}\label{eq:hinf_value}
V_{\mathrm{wc}}(\theta) = \operatorname{tr}(X W_\theta)
\end{equation}
is the robust analog of $\operatorname{tr}(PW)$, and as $\gamma\!\to\!\infty$ the
disturbance player switches off, $X\!\to\!P$ and $V_{\mathrm{wc}}\!\to\!\operatorname{tr}(PW)$,
the $H_2$ cost. Crucially, \eqref{eq:gare} is a standard DARE (only its input
weight is indefinite), so \emph{the envelope-theorem sensitivities are the very
same algebraic identities as the $H_2$ Riccati}---an identity that never used
$R\succ0$. With $S$ the game-closed-loop gramian
$S=A_{\mathrm{cl}}SA_{\mathrm{cl}}^\top+W_\theta$,
\begin{equation}\label{eq:gare_grad}
\frac{\partial V_{\mathrm{wc}}}{\partial A_\theta}=2\,X A_{\mathrm{cl}}S,\quad
\frac{\partial V_{\mathrm{wc}}}{\partial \tilde B_\theta}=2\,X A_{\mathrm{cl}}S\tilde K^\top,\quad
\frac{\partial V_{\mathrm{wc}}}{\partial W_\theta}=X,
\end{equation}
(with $\partial V_{\mathrm{wc}}/\partial Q=S$ and the $R$-block of $\tilde K S\tilde K^\top$),
chained through \eqref{eq:chain} against the model Jacobians.

The regularity of this route, unlike the minimal-$\gamma$ SDP
(Remark~\ref{rem:hinfreg}), is clean and checkable.

\begin{proposition}[Fixed-$\gamma$ $H_\infty$ regularity]\label{prop:hinfreg}
Fix $\gamma$ above the optimal attenuation and let $X\succeq0$ be the stabilizing
solution of the game GARE~\eqref{eq:gare}, with $R+B_\theta^\top XB_\theta\succ0$
and the game closed loop $A_{\mathrm{cl}}$ Schur. Then $X$ is a locally unique
$C^1$ function of $\theta$: the GARE residual $\mathcal R(X,\theta)=0$ has, at such
an $X$, a linearization in $X$ equal to the closed-loop Stein operator
$\delta X\mapsto\delta X-A_{\mathrm{cl}}^\top\delta X A_{\mathrm{cl}}$, which is
invertible precisely because $\rho(A_{\mathrm{cl}})<1$; the implicit function
theorem then gives $X(\theta)\in C^1$ and the closed-form
gradient~\eqref{eq:gare_grad} for $V_{\mathrm{wc}}(\theta)=\operatorname{tr}(XW_\theta)$.
The dual filter GARE is regular under the analogous conditions, giving
$V_{\mathrm{est}}(\theta)\in C^1$.
\end{proposition}

This is the result the MTDC study of Section~\ref{sec:sim:mtdc} relies on: it
works at a fixed $\gamma$, so the delicate minimal-$\gamma$ rank drop never enters.

The \emph{filter} GARE is the exact dual, giving the estimation half of an
output-feedback robust co-design. For the plant \eqref{eq:hinf_plant} with
measurement $y_k=C_\theta x_k+D_v v_k$ and estimation-error output $\zeta=L_e x$
($L_e=M^{1/2}$), the worst-case a-priori error covariance $\Sigma_e$ solves the
control GARE \eqref{eq:gare} on the \emph{dual} data
$(A_\theta^\top,\ [\,C_\theta^\top\ L_e^\top\,],\ W_\theta,\ \operatorname{diag}(V,-\gamma^2 I))$
--- measurement injection $C_\theta^\top$ and the estimated output $L_e^\top$
playing the roles of $B_\theta$ and $E_\theta$ --- and the worst-case estimation
cost is $V_{\mathrm{est}}(\theta)=\operatorname{tr}(M\Sigma_e)$, the robust twin of
the Kalman value $\operatorname{tr}(M\Sigma^\varepsilon)$; its gradient is the dual
of \eqref{eq:gare_grad} (contract $2\,\Sigma_e\tilde A_{\mathrm{cl}}\tilde S$ against
$\partial A_\theta^\top/\partial\theta$, and the $C_\theta,V$ blocks of
$2\,\Sigma_e\tilde A_{\mathrm{cl}}\tilde S\tilde K^\top$). This is the form used in
Section~\ref{sec:sim:mtdc}: control value $\operatorname{tr}(XW)$ and estimation
value $\operatorname{tr}(M\Sigma_e)$ with exact gradients at $O(r_x^3)$, no LMI solve
and no rank-deficient dual. The fixed-$\theta$ separation into two Riccatis is,
again, only an inner convenience: the shared $\theta$ moves $A_\theta$ and hence
both halves.

Two remarks connect this to the rest of the framework. First, the weight/noise
co-design of Section~\ref{sec:weights} carries over unchanged: the disturbance
weights $E_\theta = W_\theta^{1/2}$, $D_v = V_\theta^{1/2}$ and the performance
weights in $C_z, D_z$ are again data objects $D(\theta)$, so ``how good a sensor
to buy'' remains a continuous robust co-design variable, priced against
$J_{\mathrm{des}}$. Second, the choice between the $H_2$/LQG and $H_\infty$ inner
values is a modeling decision about the disturbance --- average-case stochastic
versus worst-case $\ell_2$ --- that ContEst treats identically at the outer level:
both expose an exact gradient from an inner dual, and the outer method of
Section~\ref{sec:algorithm} is agnostic to which is used.

%=======================================================================
\section{The locally linearized case: MPC and the eKF}\label{sec:ekf}
%=======================================================================

Beyond LQG we treat nonlinear dynamics and outputs with additive disturbances,
$x_{k+1} = \bar f(x_k,u_k;\theta_f) + w_k$, $y_k = \bar h(x_k;\theta_h) + v_k$,
with $w_k, v_k$ zero-mean of covariances $\Sigma_w, \Sigma_v$ (possibly functions
of $\theta$). The Bayesian filter
\eqref{eq:Tupdate}--\eqref{eq:Mupdate} is infinite dimensional; we approximate it
by an extended Kalman filter (eKF) \cite{anderson2012optimal,ramadan2024extended},
propagating the first two moments. With posterior mean and covariance
$x_{k\mid k}, \Sigma_{k\mid k}$, the time update is
$x_{k+1\mid k} = \bar f(x_{k\mid k},u_k)$,
$\Sigma_{k+1\mid k} = F_k \Sigma_{k\mid k} F_k^\top + \Sigma_w$, and the
measurement update is
$x_{k\mid k} = x_{k\mid k-1} + L_k(y_k - \bar h(x_{k\mid k-1}))$,
$\Sigma_{k\mid k} = \Sigma_{k\mid k-1} - L_k H_k \Sigma_{k\mid k-1}$, with
$L_k = \Sigma_{k\mid k-1} H_k^\top (H_k \Sigma_{k\mid k-1} H_k^\top +
\Sigma_v)^{-1}$ and Jacobians
$F_k = \partial \bar f/\partial x |_{x_{k\mid k}}$,
$H_k = \partial \bar h/\partial x |_{x_{k\mid k-1}}$. The eKF is exact only for
affine $\bar f, \bar h$ and Gaussian noise \cite{anderson2012optimal}.

\emph{Quadratic cost in the information state.} Take the stage cost quadratic
about a reference (set-point) trajectory $(s_k^x,s_k^u)$,
$\ell(x_k,u_k) = \|x_k - s_k^x\|_Q^2 + \|u_k - s_k^u\|_R^2$---exact for a
quadratic regulator, and otherwise the second-order Taylor model of a smooth
cost about that trajectory. Lemma~\ref{lem:stagecost} with the eKF Gaussian then
gives the following.

\begin{lem}\label{lem:ekfcost}
Under the above, $\mathbb{E}\,\ell(x_k,u_k) = \mathbb{E}\{
\|x_{k\mid k} - s_k^x\|_Q^2 + \|u_k - s_k^u\|_R^2 +
\operatorname{tr}(Q\Sigma_{k\mid k})\}$.
\end{lem}

The estimation cost thus appears as $\operatorname{tr}(Q\Sigma_{k\mid k})$,
summed over the horizon---the quantity through which the sensing parameters
$\theta_h$ act on the objective.

\emph{Convex MPC over the information state.} Stack the mean and the
upper-triangular part of the covariance into the information state
$\zeta_k = (x_{k\mid k},\, \operatorname{vech}\Sigma_{k\mid k}) \in
\mathbb{R}^{n_\zeta}$, $n_\zeta = r_x + r_x(r_x+1)/2$, with lifted dynamics
induced by the (prediction-mode) eKF map. Linearizing about the nominal
information trajectory yields $A_\theta^{\mathrm{info}}, B_\theta^{\mathrm{info}}$,
and the second-stage problem becomes the convex program---the \emph{surrogate}
$\widehat J_c$ that the outer loop actually descends, distinct from both the true
stochastic value $J_{\mathrm{stoch}}$ of \eqref{eq:Jstoch} and the realized
closed-loop cost $J_c$ evaluated by Monte-Carlo in Section~\ref{sec:numerical}
(see the exactness ledger, Table~\ref{tab:ledger}),
\begin{equation}\label{eq:mpc}
\begin{aligned}
\widehat J_c(\theta) = &\min_{\zeta_{1:N},\, u_{1:N-1}}\
\sum_{t=1}^{N-1} \ell(\zeta_t,u_t) + \ell^N(\zeta_N)\\
\text{s.t. } & \zeta_{t+1} = A_\theta^{\mathrm{info}}\zeta_t +
B_\theta^{\mathrm{info}} u_t,\quad \zeta_1 = \bar\zeta,\\
& \Sigma_{1\mid 1} \succ 0,\quad c(\zeta_{t+1},u_t) \le 0,
\end{aligned}
\end{equation}
where $\ell$ now carries the $\operatorname{tr}(Q\Sigma_{t\mid t})$ term of
Lemma~\ref{lem:ekfcost} and $c(\cdot)$ collects state/input constraints.

\emph{Design gradients from the dynamics costates.} Let $\lambda_t$ be the
multipliers (costates) of the dynamics equalities in \eqref{eq:mpc}, returned by
the solver. By Theorem~\ref{prop:envelope} the only explicit
$\theta$-dependence is through $A_\theta^{\mathrm{info}}, B_\theta^{\mathrm{info}}$,
so
\begin{equation}\label{eq:gradmpc}
\frac{\partial \widehat J_c}{\partial \theta_i} = \sum_t \lambda_t^\top \Big(
\frac{\partial A_\theta^{\mathrm{info}}}{\partial \theta_i}\zeta_t^\star +
\frac{\partial B_\theta^{\mathrm{info}}}{\partial \theta_i} u_t^\star \Big).
\end{equation}
The factors $\partial A_\theta^{\mathrm{info}}/\partial\theta_i$ and
$\partial B_\theta^{\mathrm{info}}/\partial\theta_i$ are Jacobians of the
eKF-induced linearization, obtained by automatic differentiation of the
\emph{filter map}, never of the optimizer.

\begin{rem}
Equations \eqref{eq:gradAB}--\eqref{eq:chain} and \eqref{eq:gradmpc} are one
principle in two guises: the inner dual (LMI dual $S$, or costate $\lambda$)
contracts against a model Jacobian to give the exact design gradient. This is the
computational heart of ContEst.
\end{rem}

\begin{rem}[Where exactness stops: the nonlinear case]\label{rem:ekfbreaks}
Theorem~\ref{prop:envelope} applies verbatim to the LQG SDPs, which are genuinely
convex; in the nonlinear regime it applies to the convex program \eqref{eq:mpc},
not to the underlying stochastic control value \eqref{eq:Jstoch}. Three gaps
separate the two, and we are explicit about them. \emph{(i)~Filtering
approximation:} the eKF replaces the exact Bayesian recursion
\eqref{eq:Tupdate}--\eqref{eq:Mupdate} by its first two moments, exact only for
affine $\bar f,\bar h$ and Gaussian noise \cite{anderson2012optimal}.
\emph{(ii)~Non-separability and nonconvexity:} outside LQG the separation of
Section~\ref{sec:lqg} fails, and $J_{\mathrm{stoch}}(\theta)$ is in general
nonconvex and possibly nonsmooth in $\theta$, so a single stationary point need
not be global---hence the multi-start outer loop of
Section~\ref{sec:numerical}. \emph{(iii)~Linearization:} $A_\theta^{\mathrm{info}},
B_\theta^{\mathrm{info}}$ are taken about a nominal information trajectory, so
\eqref{eq:gradmpc} is the \emph{exact} gradient of the convex surrogate but an
\emph{approximate} gradient of the true value; the surrogate is what the outer
method actually descends. Assumption~\ref{as:env} is therefore imposed on
\eqref{eq:mpc}, and the gradient check of Section~\ref{sec:numerical} confirms it
holds for the surrogate. Making the certainty-equivalent
and linearization gaps quantitative---error bounds relating the surrogate
stationary point to a stationary point of \eqref{eq:Jstoch}---is left as future
work; a sharper information state than the eKF (e.g.\ the eKF--Koopman lift
\cite{ramadan2024extended}) narrows gap~(i).
\end{rem}

Table~\ref{tab:ledger} is the resulting \emph{exactness ledger}: it separates the
three objects the analysis involves---the true stochastic value
$J_{\mathrm{stoch}}$ of \eqref{eq:Jstoch}, the convex surrogate $\widehat J_c$ of
\eqref{eq:mpc}, and the realized closed-loop cost $J_c$ that
Section~\ref{sec:numerical} evaluates by Monte-Carlo---and states, per regime, what
the envelope gradient \eqref{eq:envelope}/\eqref{eq:gradmpc} is \emph{exact} for
and what gap remains to the true value. In the LQG and fixed-$\gamma$ $H_\infty$
regimes the inner value is the true stationary cost and the gradient is exact
(Propositions~\ref{prop:lqgreg},~\ref{prop:hinfreg}); in the nonlinear regime the
gradient is exact for the surrogate $\widehat J_c$ only, and the Monte-Carlo
evaluation of Section~\ref{sec:numerical} does theoretical work by empirically
bounding the surrogate gap---each optimized design is shown to improve the
\emph{realized} closed-loop cost, not merely the surrogate.

\begin{table}[t]
\centering
\caption{Exactness ledger. $J_{\mathrm{stoch}}$: true stochastic value
\eqref{eq:Jstoch}; $\widehat J_c$: convex surrogate \eqref{eq:mpc}; $J_c$:
realized Monte-Carlo closed-loop cost (Section~\ref{sec:numerical}). The envelope
gradient is exact for the ``exact for'' column; the last column names the residual
gap to $J_{\mathrm{stoch}}$.}
\label{tab:ledger}
\small
\begin{tabular}{@{}p{1.75cm}p{1.55cm}p{1.55cm}p{1.9cm}@{}}
\toprule
regime & inner problem & gradient exact for & gap to true value\\
\midrule
LQG (Sec.~\ref{sec:lqg}) & SDPs/Riccati & $J_{\mathrm{cont}},J_{\mathrm{est}}$ = true LQG cost & none\\
$H_\infty$ fixed $\gamma$ (Sec.~\ref{sec:hinf}) & game GARE & $V_{\mathrm{wc}},V_{\mathrm{est}}$ & none (given $\ell_2$ model)\\
nonlinear (Sec.~\ref{sec:ekf}) & convex prog.~\eqref{eq:mpc} & surrogate $\widehat J_c$ & eKF closure + linearization + CE policy\\
\bottomrule
\end{tabular}
\end{table}

%=======================================================================
\section{Sensor allocation: discrete selection and sparse relaxation}
\label{sec:sensor}
%=======================================================================

A recurring design decision is \emph{which} sensors to deploy. With $r_y$
candidate sensors, binary indicators $b_j\in\{0,1\}$ (sensor $j$ contributing its
row $C_j$ and noise variance $\sigma_{v,j}^2$ when $b_j=1$) and per-sensor costs
$c_j$ make the co-design $\min_{\theta,b\in\{0,1\}^{r_y}}
J_{\mathrm{stoch}}(\theta,b)+c^\top b$ a mixed-integer (SDP) program, whose exact
solution is combinatorial (NP-hard, $2^{r_y}$ subsets)
\cite{joshi2009sensor,tzoumas2016sensor}. Standard tractable remedies---greedy
selection with submodularity guarantees \cite{tzoumas2016sensor,zhang2017sensor},
the convex Boolean relaxation $b_j\in[0,1]$ \cite{joshi2009sensor}, and
sparsity-promoting $\ell_1$ \cite{dhingra2014admm,saraf2017h2,das2017sparse}---are
all compatible with ContEst; we adopt the $\ell_1$ route, which fits it with no
change to the gradient machinery.

Replace the indicators by continuous gains $\alpha_j\ge0$ scaling each sensor's
contribution, $C(\alpha)=\operatorname{diag}(\alpha)\,C_{\mathrm{base}}$
($\alpha_j=0$ removing sensor $j$), so $\alpha$ joins $\theta$ through $C_\theta$,
and put the sparsity penalty in the \emph{design cost},
\begin{equation}\label{eq:Jdes}
J_{\mathrm{des}}(\theta) = J_{\mathrm{hw}}(\theta) + \lambda_s \| \alpha \|_1,
\qquad \alpha \ge 0,
\end{equation}
with $J_{\mathrm{hw}}$ the smooth design penalties and $\lambda_s>0$ trading
estimation quality against the number of active sensors. Because the $\ell_1$ term
sits in $J_{\mathrm{des}}$ and not the inner problem, the envelope gradients
\eqref{eq:chain}, \eqref{eq:gradmpc} are untouched. Its one kink is removed by the
standard epigraph lift \cite{boyd2004convex}: substitute
$\lambda_s\|\alpha\|_1\mapsto\lambda_s\mathbf 1^\top t$ and append
$\alpha\preceq t$ ($\alpha\ge0$), which is smooth (linear in the slacks $t$) over a
polyhedron, so the (L-)BFGS of Section~\ref{sec:algorithm} keeps its convergence
guarantees; equivalently one takes a proximal (soft-threshold) step on
$\lambda_s\|\alpha\|_1$. When a binary deployment is required, the converged
$\alpha^\star$ is thresholded (optionally with a fixed-sensor re-solve, or
reweighted-$\ell_1$ iterations \cite{candes2008reweighted}). ContEst thus treats
sensor allocation and continuous sensor tuning uniformly; the MTDC study of
Section~\ref{sec:sim:mtdc} exercises this allocation on the robust $H_\infty$ path,
pruning $60\%$ of the DC-voltage sensors at no estimation cost.

%=======================================================================
\section{Outer optimization}\label{sec:algorithm}
%=======================================================================

Problem~\ref{prob:bilevel} is solved by a projected quasi-Newton method on
$\theta$ over the box $\Theta = [\theta_{\mathrm{lb}},\theta_{\mathrm{ub}}]$. Each
outer iteration: (1)~solves the inner convex problem
(\eqref{eq:Jcont}/\eqref{eq:Jest} or \eqref{eq:mpc})---one solve, or one per
scenario when several are averaged (Remark~\ref{rem:scenarios});
(2)~reads off its duals ($S$ or $\lambda$); (3)~assembles the design
gradient via \eqref{eq:chain} or \eqref{eq:gradmpc}; (4)~adds
$\nabla J_{\mathrm{des}}$, including the $\ell_1$ subgradient of
\eqref{eq:Jdes}; and (5)~takes a quasi-Newton (e.g.\ L-BFGS) step with projection
onto $\Theta$. Because the gradients are exact---neither finite-difference nor
approximate implicit derivatives---the outer iteration retains the fast local
convergence of quasi-Newton methods, and the dominant per-iteration cost is the
inner convex solves, which scale polynomially.

\emph{Stationarity.} When the objective is smooth---the LQG and fixed-$\gamma$
$H_\infty$ regimes, where Propositions~\ref{prop:lqgreg} and~\ref{prop:hinfreg}
make $J\in C^1$, and the nonlinear surrogate~\eqref{eq:mpc} likewise---every
accumulation point of the projected quasi-Newton iterates is a stationary point of
$J$ over $\Theta$ \cite{beck2017first}. The one nonsmooth term, the sparsity
penalty $\lambda_s\|\alpha\|_1$ of \eqref{eq:Jdes}, is handled without forfeiting
this guarantee by the epigraph lift of Section~\ref{sec:sensor} (equivalently, a
proximal-gradient step on the composite objective \cite{beck2017first}).

%=======================================================================
\section{Simulation studies}\label{sec:numerical}
%=======================================================================

We evaluate ContEst on three nonlinear co-design problems spanning very different
physics: spacecraft attitude determination and control (Section~\ref{sec:sim:adcs}), a
binary distillation column (Section~\ref{sec:sim:distill}), and PLL/DSE estimator
tuning in a low-inertia grid (Section~\ref{sec:sim:pll}), followed by a
\emph{large-scale, linear} power-system study on the steady-state \emph{robust}
($H_\infty$) inner problem (multi-terminal HVDC droop coordination,
Section~\ref{sec:sim:mtdc}) and a complexity analysis of the two inner
problems under interior-point solvers (Section~\ref{sec:sim:complexity}). Each of
these three is genuinely nonlinear, so all
use the eKF--MPC second stage of Section~\ref{sec:ekf}: the convex program
\eqref{eq:mpc} is solved with the Clarabel interior-point solver, and the design
gradient is assembled from the dynamics costates through \eqref{eq:gradmpc}.
The outer loop is the projected quasi-Newton method of
Section~\ref{sec:algorithm} (L-BFGS) over the box $\Theta$, run from $5$ initial
points---the baseline design $\theta_{\mathrm{nom}}$ plus four uniform random
restarts in $\Theta$---keeping the best minimizer. For the ADCS and PLL
studies all five starts \emph{converge to the same minimizer}, so the reported
design is the global optimum over $\Theta$ to numerical tolerance; the distillation
study is genuinely multimodal (the restarts resolve three distinct minima), and we
report the best, the restarts escaping the local optimum the baseline start alone
would return. The envelope-theorem gradients are exact, so we do not verify them
per study; as a one-time sanity check we confirmed, across all studies below, that
the analytic gradient matches central finite differences to a relative error below
$10^{-4}$ (best case $1.6\times10^{-6}$). Code reproducing every study is publicly
available.\footnote{\url{https://github.com/anonymous/ContEst} (placeholder
repository; the final URL will be provided on acceptance).}

In every study the design splits as $\theta=(\theta_f,\theta_h)$, where
$\theta_f$ enters the dynamics $\bar f$ (actuation / plant tuning, the
classical CCD axis) and $\theta_h$ is the estimation axis that makes the problem
ContEst. Each axis is posed so that its optimum is a genuine \emph{interior} trade
---a scarce resource allocated under a budget, or a saturating benefit balanced
against a provisioning cost---rather than a monotone ``more is better'' knob pinned
at a bound. The reported costs are \emph{Monte-Carlo averages of the realized
closed-loop cost}: for each design we draw independent initial states and
process/measurement noise, run the true nonlinear plant forward under a
receding-horizon controller with the eKF processing the noisy measurements, and
average the incurred cost over the sampled trajectories. We decompose it into its
two physical parts, both evaluated on the \emph{true} state $x_t$,
\begin{equation}\label{eq:costsplit}
\begin{aligned}
J_{\mathrm{det}} &= \sum_{t}\!\big(\|x_t-s^x_t\|_Q^2+\|u_t-s^u_t\|_R^2\big),\\
J_{\mathrm{est}} &= \sum_{t}\|x_t-\hat x_{t\mid t}\|_Q^2,
\end{aligned}
\end{equation}
the (deterministic) control cost and the realized estimation error against the eKF
mean $\hat x_{t\mid t}$ (whose expectation is $\sum_t\operatorname{tr}(Q\Sigma_{t\mid t})$),
with $J_c=J_{\mathrm{det}}+J_{\mathrm{est}}$ and reported total
$J_{\mathrm{tot}} = J_c + J_{\mathrm{des}}$ (single scenario, $M=1$;
Remark~\ref{rem:scenarios}). The \emph{baseline} design is driven by a
certainty-equivalence controller that uses \emph{only the eKF mean} (it plans on
the mean dynamics and is blind to the covariance), whereas the ContEst design is
driven by the covariance-aware information-state MPC of Section~\ref{sec:ekf}; the
reported reduction therefore reflects both the re-designed parameters and the
estimation-aware control they enable. (The distillation study, whose many-stage
information-state MPC is expensive to simulate in closed loop, instead reports the
open-loop surrogate $\widehat J_c$ of \eqref{eq:mpc} directly, with
$J_{\mathrm{est}}=\sum_t\operatorname{tr}(Q\Sigma_{t\mid t})$ from the eKF
covariance rollout.) The design optimization itself is unchanged---multi-start
L-BFGS on the envelope-gradient objective. As anticipated by the exactness ledger
(Table~\ref{tab:ledger}), the surrogate $\widehat J_c$ is an \emph{open-loop}
planned cost and is therefore optimistic relative to the noise-driven realized cost
$J_c$; what matters for design is that minimizing the surrogate also lowers the realized
cost, which the Monte-Carlo evaluation confirms in every study, with the design
ranking preserved.

%-----------------------------------------------------------------------
\subsection{Spacecraft attitude determination and control (ADCS)}
\label{sec:sim:adcs}
%-----------------------------------------------------------------------

\emph{Significance.} Pointing accuracy on every space mission is ultimately set
by both the attitude \emph{estimator} (the sensor suite: star trackers, gyros)
and the reaction-wheel authority. On a spacecraft these compete for a shared,
tightly constrained \emph{power/mass budget}: a heavier, higher-power star tracker
or gyro reports lower noise, and a larger wheel delivers more torque, but every
watt and kilogram spent on one is unavailable to the others. The benefit of each
saturates, so the design question is not ``how good a sensor'' but how to
\emph{split one budget} across sensing and actuation---an allocation classical
control co-design cannot pose because it has no sensing axis.

\emph{Formulation.} A rigid satellite with asymmetric inertia
$J=\operatorname{diag}(4,6,5)\,\si{kg.m^2}$ regulates attitude and rate to zero.
The state $x=[\phi,\theta_a,\psi,\omega_x,\omega_y,\omega_z]$ ($r_x=6$) uses
small-angle attitude and body rates; the input $u$ is three reaction-wheel
torques ($r_u=3$); the outputs are a star tracker (attitude) and a rate gyro
($r_y=6$). The genuine nonlinearity is the gyroscopic coupling $\omega\times
(J\omega)$, which vanishes at the origin so $\bar f(0,0;\theta)=0$ for all
$\theta$; integration is at $\SI{0.1}{s}$ over $N=15$. The three design
parameters---reaction-wheel authority $e_{\mathrm{rw}}$ (in $\bar f$) and the
star-tracker and rate-gyro precisions $\alpha_{\mathrm{st}},\alpha_{\mathrm{gyro}}$
(entering the measurement covariance $V(\theta)$)---draw on a single resource
budget. Attitude error is weighted four times more heavily than rate, so the star
tracker is intrinsically more valuable than the gyro.

\emph{State-space model.} With $\Delta t=\SI{0.1}{s}$ the system~\eqref{eq:ss}
reads $x_{k+1}=\bar f(x_k,u_k;\theta)+w_k$, $y_k=\bar h(x_k)+v_k$ with
\begin{subequations}\label{eq:adcs-model}
\begin{align}
(\phi,\theta_a,\psi)^{+} &= (\phi,\theta_a,\psi) + \Delta t\,(\omega_x,\omega_y,\omega_z),\\
\omega^{+} &= \omega + \Delta t\,J^{-1}\!\big(e_{\mathrm{rw}}\,u - \omega\times(J\omega)\big),\\
\bar h &= \big[(\phi,\theta_a,\psi),\ (\omega_x,\omega_y,\omega_z)\big],
\end{align}
\end{subequations}
the gyroscopic term $\omega\times(J\omega)$ the nonlinearity (zero at
$\omega=0$), and $\theta$-dependent sensor noise
$V(\theta)=\operatorname{diag}(v_0/\alpha_{\mathrm{st}}^2\,I_3,\,v_0/\alpha_{\mathrm{gyro}}^2\,I_3)$
(better precision $\Rightarrow$ lower variance, saturating). \emph{Noises:}
$w_k\sim\mathcal N(0,\Sigma_w)$ with
$\Sigma_w=\operatorname{diag}(10^{-6}I_3,10^{-4}I_3)$, base sensor variance
$v_0=0.20$, prior $\Sigma_0=2\!\times\!10^{-2}I_6$. \emph{Cost:}
$Q=\operatorname{diag}(200I_3,50I_3)$, $R=I_3$, with $|\tau_i|\le\SI{0.8}{N.m}$.
\emph{Design cost:} a shared budget $c_b(e_{\mathrm{rw}}+\alpha_{\mathrm{st}}+
\alpha_{\mathrm{gyro}}-B)^2$ ($c_b=30$, $B=3$), with box
$e_{\mathrm{rw}}\in[0.3,3]$, $\alpha_{\mathrm{st}},\alpha_{\mathrm{gyro}}\in[0.3,5]$,
and baseline $\theta_{\mathrm{nom}}=\mathbf 1$ (budget spent uniformly).

\emph{Results.} All five starts agreed.
The co-design (Table~\ref{tab:adcs}) reallocates the budget:
it upgrades the star tracker most ($\alpha_{\mathrm{st}}\!:\,1\to1.64$), trims the
lightly weighted rate gyro ($\alpha_{\mathrm{gyro}}\!:\,1\to0.78$), and modestly
raises wheel authority ($e_{\mathrm{rw}}\!:\,1\to1.19$)---none pinned at a bound.
This cuts the estimation cost $J_{\mathrm{est}}$ by $23.8\%$ ($173.1\to132.0$) and
the control cost $J_{\mathrm{det}}$ by $4.0\%$ ($448.3\to430.5$), for a $9.5\%$
reduction in $J_c=J_{\mathrm{det}}+J_{\mathrm{est}}$ ($621.4\to562.5$) and a $7.7\%$
reduction in the total $J_{\mathrm{tot}}$ ($621.4\to573.5$).

\emph{Comments.} With a single shared budget the axes genuinely compete, and the
optimizer spends it where the cost weights make it count---attitude sensing (the
star tracker) over rate sensing (the gyro), with a little left for wheel
authority. The interior split is exactly the effort allocation \emph{across} the
sensing and actuation axes that a separate, control-only design cannot represent.
A sensor-selection method \cite{joshi2009sensor,tzoumas2016sensor} cannot pose
this trade because it has no actuation axis to spend the shared budget against;
plant--controller co-design cannot pose it because it has no sensing axis at all.

\begin{table}[t]
\centering
\caption{Spacecraft ADCS: baseline (uniform budget) vs.\ co-designed allocation of
the shared power/mass budget $e_{\mathrm{rw}}+\alpha_{\mathrm{st}}+\alpha_{\mathrm{gyro}}\approx B=3$.}
\label{tab:adcs}
\begin{tabular}{llcc}
\hline
parameter & axis & baseline & optimal \\
\hline
$e_{\mathrm{rw}}$       & $f$ (reaction wheels) & 1.000 & 1.186 \\
$\alpha_{\mathrm{st}}$  & $V$ (star tracker)    & 1.000 & 1.642 \\
$\alpha_{\mathrm{gyro}}$& $V$ (rate gyro)       & 1.000 & 0.777 \\
\hline
\end{tabular}
\end{table}

%-----------------------------------------------------------------------
\subsection{Binary distillation column}
\label{sec:sim:distill}
%-----------------------------------------------------------------------

\emph{Significance.} Distillation is the workhorse separation of the process
industries and its control is dominated by two classical layout decisions: on
\emph{which} tray to introduce the feed, and on \emph{which} tray to place the
inferential (temperature) sensor that stands in for an unmeasured composition
profile. Both are \emph{placement} choices with well-known interior optima---a
feed stage too high or too low wastes separation, and the most informative
temperature tray is the steep mid-column stage, not an end tray. The feed
locations shape the plant $\bar f$ while the sensor locations shape the output map
$\bar h$, so this is a ContEst problem in which \emph{neither} axis is a
gain: every parameter moves a discrete-in-spirit location on a continuous scale.

\emph{Formulation.} A twelve-stage column (stage~1 reboiler
$\dots$ stage~12 condenser) separates a light/heavy mixture; the state
$x=[\Delta x_1,\dots,\Delta x_{12}]$ ($r_x=12$) collects stage light-component
mole-fraction deviations about the nominal profile, in deviation form so
$\bar f(0,0;\theta)=0$ for every $\theta$. The nonlinearity is genuine and
enters twice: the vapour--liquid equilibrium
$y(x)=\alpha x/(1+(\alpha-1)x)$ ($\alpha=2.5$) and the bubble-point temperature
$T(x)=T_{B0}-\Delta T_B\,y(x)$. The single input is the reboiler boil-up
deviation $u=\Delta V$ ($r_u=1$), limited to $|\Delta V|\le1.5$; the measurements
($r_y=3$) are three Gaussian-weighted stage temperatures read at the sensor
locations. Integration is at $\SI{0.05}{s}$ over $N=12$. The design carries
$N_f=3$ feed-stage locations $\theta_f=(f_1,f_2,f_3)$ and $N_s=3$
temperature-sensor locations $\theta_h=(s_1,s_2,s_3)$, all in $[2,11]$; the total
feed is distributed over trays by a \emph{sum} of normalised Gaussians centred at
the $f_j$, so the cumulative feed above each stage sets that stage's internal
liquid traffic $L_i(\theta_f)$ and the feed placement nonlinearly reshapes the
whole convective coupling of the column. The upset is a feed-composition
disturbance that peaks mid-column, and the stage weights $Q$ are heaviest there.
Because temperature is the \emph{only} measurement, where the three sensors sit
sets how well the entire composition profile is reconstructed, and they must spread
to cover it rather than pile onto one tray. Placement is nearly free to relocate,
so $J_{\mathrm{des}}$ carries only a mild regularisation
$\lambda_{\mathrm{des}}\lVert\theta-\theta_{\mathrm{nom}}\rVert^2$ about the naive
default $\theta_{\mathrm{nom}}$ with feeds at $(3,4.5,6)$ and sensors at
$(2,3.5,5)$ (feeds below centre, sensors low near the bottoms product)---the win is
driven by the column physics, not by the design penalty.

\emph{State-space model.} With equilibrium and temperature maps
$y(x)=\alpha x/(1+(\alpha-1)x)$, $T(x)=T_{B0}-\Delta T_B\,y(x)$, equilibrium
deviations $\Delta y_i=y(x_i^{\mathrm{ss}}+\Delta x_i)-y_i^{\mathrm{ss}}$ and
$\Delta t=\SI{0.05}{s}$, the system~\eqref{eq:ss} reads
$x_{k+1}=\bar f(x_k,u_k;\theta)+w_k$, $y_k=\bar h(x_k;\theta)+v_k$ with, for
$i=1,\dots,12$ and $j=1,\dots,3$,
\begin{subequations}\label{eq:distill-model}
\begin{align}
\Delta x_i^{+} &= \Delta x_i + \tfrac{\Delta t}{M}\big(\mathrm{liq}_i+\mathrm{vap}_i\big),\\
\mathrm{liq}_i &= L_i(\theta_f)\,(\Delta x_{i+1}-\Delta x_i),\\
\mathrm{vap}_i &= (V+u)\,(\Delta y_{i-1}-\Delta y_i) + u\,(y_{i-1}^{\mathrm{ss}}-y_i^{\mathrm{ss}}),\\
\bar h_j &= \textstyle\sum_{i=1}^{12}\hat w_i(s_j)\,\big(T(x_i^{\mathrm{ss}}+\Delta x_i)-T_i^{\mathrm{ss}}\big),
\end{align}
\end{subequations}
boundary terms $\Delta x_{13}=\Delta y_0=y_0^{\mathrm{ss}}=0$, feed weights
$w_k(\theta_f)\propto \sum_{j=1}^{3}e^{-(k-f_j)^2/2\sigma_f^2}$ (normalized)
setting the internal liquid traffic
$L_i(\theta_f)=L_{\mathrm{ref}}+F\sum_{k>i}w_k$, and sensor weights
$\hat w_i(s_j)\propto e^{-(i-s_j)^2/2\sigma_{\mathrm{sens}}^2}$
(normalized)---the maps carrying the design coupling. Constants $\alpha=2.5$,
$L_{\mathrm{ref}}=2$, $V=2.5$, $F=2.5$, $M=1$,
$\sigma_f=\sigma_{\mathrm{sens}}=1.2$, $T_{B0}=\SI{380}{K}$,
$\Delta T_B=\SI{40}{K}$. \emph{Noises:} $\Sigma_w=2\!\times\!10^{-5}I_{12}$,
$\Sigma_v=0.6\,I_3$ ($\si{K^2}$), prior $\Sigma_0=10^{-2}I_{12}$. \emph{Cost:}
$Q_{ii}=20+30\,e^{-(i-6.5)^2/2(12/5)^2}$ (heaviest mid-column), $R=2$,
$|u|\le1.5$, with a mild placement regularization
$\lambda_{\mathrm{des}}\lVert\theta-\theta_{\mathrm{nom}}\rVert^2$.

\emph{Results.} The design landscape is multimodal: the five-start L-BFGS resolves three distinct
minima, and the baseline start alone converges to an inferior one, so the random
restarts are what recover the reported optimum. ContEst spreads the feed
streams---two settle lower toward the reboiler ($f_{1,2}:3.0,4.5\to2.5,3.6$) while
the third climbs high up the column ($f_3:6.0\to9.0$)---and fans the three
temperature sensors out of the bottoms region across the steep mid-and-upper trays
($s:(2.0,3.5,5.0)\to(3.0,5.4,8.0)$), covering the composition profile rather than
clustering low (Table~\ref{tab:distill}). This cuts the estimation cost
$J_{\mathrm{est}}$ by $22.8\%$ ($12.8\to9.9$) and the control cost
$J_{\mathrm{det}}$ by $1.9\%$, for a $10.2\%$ reduction in $J_c$ and---after the
small placement penalty---a $6.5\%$ reduction in the total $J_{\mathrm{tot}}$
($32.3\to30.2$).

\emph{Comments.} This study stresses ContEst on the \emph{placement} form of the
sensing axis---three movable inferential sensors and three feed streams---on a
nonconvex, multimodal objective. Both physical effects that make distillation
layout hard---the interior-optimal feed distribution and the interior-optimal
sensor trays---are recovered jointly by the same envelope gradient, and the gain is
estimation-led ($J_{\mathrm{est}}$ down $22.8\%$ against a nearly flat
$J_{\mathrm{det}}$), reflecting that with temperature the only measurement the
reconstruction quality dominates the achievable regulation. It also illustrates why
co-design matters even without a hard sensor budget: a control-only design would
inherit the naive low sensor trays and forfeit the estimation-cost reduction, while
the baseline start alone would stall at a local minimum. A discrete
sensor-selection method cannot represent this design either, because the feed and
sensor trays are \emph{continuous} locations whose interior optima fall between
integer stages---the coupling of $\bar f$ and $\bar h$ through the column physics is
invisible to menu enumeration.

\begin{table}[t]
\centering
\caption{Binary distillation: naive baseline vs.\ co-designed placements. All six
parameters are continuous tray locations (not gains); baseline $\theta_{\mathrm{nom}}$
places feeds at $(3,4.5,6)$ and sensors at $(2,3.5,5)$.}
\label{tab:distill}
\begin{tabular}{llcc}
\hline
parameter & axis & baseline & optimal \\
\hline
$f_1$ & $f$ (feed stage)    & 3.000 & 2.479 \\
$f_2$ & $f$ (feed stage)    & 4.500 & 3.584 \\
$f_3$ & $f$ (feed stage)    & 6.000 & 9.018 \\
$s_1$ & $h$ (temp.\ sensor) & 2.000 & 2.972 \\
$s_2$ & $h$ (temp.\ sensor) & 3.500 & 5.428 \\
$s_3$ & $h$ (temp.\ sensor) & 5.000 & 8.010 \\
\hline
\end{tabular}
\end{table}

%-----------------------------------------------------------------------
\subsection{Estimator co-design: PLL/DSE tuning in a low-inertia grid}
\label{sec:sim:pll}
%-----------------------------------------------------------------------

\emph{Significance.} The studies above move the sensing axis $\theta_h$ through
the output map---sensor precision or placement, i.e.\ \emph{what} or \emph{where}
one senses. This study co-designs the \emph{estimator itself}---the per-inverter
dynamic state estimation (DSE). Every
grid-following (GFL) inverter estimates grid frequency and phase with a
phase-locked loop (PLL), whose bandwidth is a classic bias--variance knob: a
narrow-band PLL is smooth but sluggish and lags real frequency excursions, while
a wide-band PLL tracks fast but passes more measurement noise into the frequency
estimate---and in low-inertia, inverter-dominated grids this tuning is decisive
for stability and for fast frequency response \cite{zhao2019power}. The PLL
bandwidth is therefore an \emph{estimator} design variable that classical control
co-design cannot represent and sensor placement cannot capture. It is also the
first study to exercise the noise-level extension of
Section~\ref{sec:weights}: the PLL bandwidth enters both the dynamics
$\bar f$ (the PLL tracking state) \emph{and} the measurement covariance
$V(\theta)$ (the frequency-readout noise), so $\theta_h$ shapes $V$ rather than
$\bar h$.

\emph{Formulation.} Three GFL-inverter buses form a low-inertia chain. The state
$x=[\Delta\delta_{1:3},\Delta\omega_{1:3},\Delta\hat\omega^{\mathrm p}_{1:3}]$
($r_x=9$) stacks bus angle deviations, frequency deviations, and the
\emph{per-bus PLL frequency estimate} $\Delta\hat\omega^{\mathrm p}_i$; the input
$u=[\Delta P_{1:3}]$ is the inverter damping power ($r_u=3$). Each bus is measured
by a phasor measurement unit (PMU) angle channel and its PLL frequency channel
($r_y=6$). The genuine
nonlinearity is the synchronizing power $\propto\sin(\delta^{\mathrm{eq}}_{ij}+
\Delta\delta_i-\Delta\delta_j)$, in deviation form so $\bar f(0,0;\theta)=0$ for
all $\theta$; integration is at $\SI{0.05}{s}$ over $N=12$. The design parameters
are the per-bus inverter damping effectiveness
$\theta_f=(e_1,e_2,e_3)\in[0.3,3]^3$ and the per-bus PLL bandwidths
$\theta_h=(b_1,b_2,b_3)\in[2,20]\,\si{rad/s}$. The disturbance is concentrated at
bus~1 (large post-fault kick and prior), while buses~2,3 are quiet, so the design
must trade tracking speed against readout noise \emph{per bus}.

\emph{State-space model.} With $\Delta t=\SI{0.05}{s}$ the system reads
$x_{k+1}=\bar f(x_k,u_k;\theta)+w_k$, $y_k=\bar h(x_k)+v_k$ with, for $i=1,2,3$,
\begin{subequations}\label{eq:pll-model}
\begin{align}
\Delta\delta_i^{+} &= \Delta\delta_i + \Delta t\,\Delta\omega_i,\\
\Delta\omega_i^{+} &= \Delta\omega_i + \tfrac{\Delta t}{2H_i}
   \big(\!-P_i^{\mathrm e}(\Delta\delta) - D_i\Delta\omega_i + e_i u_i\big),\\
\Delta\hat\omega^{\mathrm p,+}_i &= \Delta\hat\omega^{\mathrm p}_i
   + \Delta t\,b_i\,(\Delta\omega_i - \Delta\hat\omega^{\mathrm p}_i),\\
\bar h &= \big[\,\Delta\delta_i;\ \Delta\hat\omega^{\mathrm p}_i\,\big]_{i=1}^{3},
\end{align}
\end{subequations}
where $P_i^{\mathrm e}(\Delta\delta)=\sum_{j\neq i}K_{ij}\big[\sin(\Delta^{\mathrm{eq}}_{ij}
+\Delta\delta_i-\Delta\delta_j)-\sin\Delta^{\mathrm{eq}}_{ij}\big]$ is the synchronizing
power ($H_i,D_i,\delta^{\mathrm{eq}}_i$ the machine inertias, dampings and
equilibrium angles) and the third line is the first-order PLL tracking the true
frequency with bandwidth $b_i$. \emph{Noise:} $\Sigma_w=\operatorname{blkdiag}
(2\!\times\!10^{-5}I_3,5\!\times\!10^{-4}I_3,5\!\times\!10^{-4}I_3)$; prior
$\Sigma_0$ large at bus~1 ($3\!\times\!10^{-2}$), small elsewhere
($6\!\times\!10^{-3}$). \emph{Cost:} $Q=\operatorname{blkdiag}(15I_3,120I_3,0_3)$
(weighting the \emph{true} angle and, heavily, the true frequency; PLL states
unweighted), $R=I_3$, $|u_i|\le1.5$. \emph{Design cost:}
$J_{\mathrm{des}}=\lambda_e\sum_i(e_i-1)^2+\lambda_b\sum_i(b_i-b_{\mathrm{nom}})^2$,
$\lambda_e=0.5$, $\lambda_b=0.02$, $b_{\mathrm{nom}}=6$. The design coupling is
that the PLL bandwidth $b_i$ enters \emph{both} $\bar f$ (the tracking line above)
and the measurement covariance, which is \emph{design-dependent},
\begin{equation}\label{eq:pll-V}
V(\theta)=\operatorname{blkdiag}\!\big(\sigma_\delta^2 I_3,\
\operatorname{diag}_i\,\sigma_\omega^2(1+\kappa\,b_i)\big),
\end{equation}
$\sigma_\delta^2=4\!\times\!10^{-3}$, $\sigma_\omega^2=2\!\times\!10^{-3}$,
$\kappa=0.30$: a wider-band PLL tracks faster but yields a noisier frequency
readout---the bias--variance knob $\theta_h$ resolves.

\emph{Results.} Here the envelope-theorem gradient flows through the
\emph{design-dependent covariance} $V(\theta)$---the numerical exercise of the
noise-level extension of Section~\ref{sec:weights}. All five starts converged to the same minimizer, so the
reported design is the global optimum over $\Theta$ to numerical tolerance. The
naive uniform PLL ($b=6$) is too sluggish; ContEst raises every bandwidth but
tunes each to its local signal-to-noise, most aggressively at the disturbed bus
($b_1\!:\,6\to13.7$) and moderately at the quiet buses
($b_{2,3}\!:\,6\to10.6$)---and \emph{none} reaches the upper bound of $20$,
because beyond the optimum the added readout noise in $V(\theta)$ outweighs faster
tracking (Table~\ref{tab:pll}). It simultaneously raises the
inverter damping authority ($e_1\to3.0$, $e_{2,3}\to1.95,2.19$). This cuts the
estimation cost $J_{\mathrm{est}}$ by $20.9\%$ ($26.6\to21.0$) and the control cost
$J_{\mathrm{det}}$ by $40.6\%$ ($151.9\to90.2$), for a $37.7\%$ reduction in
$J_c=J_{\mathrm{det}}+J_{\mathrm{est}}$ ($178.5\to111.2$) and a $34.8\%$ reduction in
the total $J_{\mathrm{tot}}$ ($178.5\to116.4$).

\emph{Comments.} This is the most vivid demonstration that the \emph{estimator} is
a first-class design variable in ContEst. The tuned quantity is neither a sensor
location nor a static gate but the bandwidth of a \emph{dynamic} estimator
embedded in every inverter, and it acts on the objective through the measurement
covariance $V(\theta)$---a channel absent from every prior study and from
control-only co-design. That the optimal bandwidths sit strictly inside the box is
the signature of the bias--variance trade-off the framework now resolves
automatically, allocating the most estimator bandwidth exactly where the dynamics
are fastest and least certain. No sensor-placement method can tune this variable,
because it tunes an estimator \emph{state} in $\bar f$ and the noise level
$V(\theta)$ rather than \emph{which} sensor to read.

\emph{Beyond the naive baseline: a competent sequential pipeline.} A practitioner
would not use the uniform baseline but a \emph{sequential} pipeline---tuning the
estimator $\theta_h$ first for a pure estimation criterion (the eKF covariance
rollout $\sum_t\operatorname{tr}(Q\Sigma_{t\mid t})$) and then the control $\theta_f$
with $\theta_h$ frozen, optionally \emph{alternating} the two. We ran both against the
joint optimum; here the designs nearly coincide (the joint optimum improves the
surrogate it descends by only $0.3\%$ over the sequential design and ties it within
Monte-Carlo noise on the realized cost). This study is, in practice, close to
\emph{separable}: the estimation-first tuning already sits near the joint optimum. Two
points follow. First, the large reductions reported above are against the \emph{naive}
design that deployed systems actually carry, the relevant baseline. Second---the
structural advantage---ContEst does not need to \emph{know} in advance whether a
problem separates: one joint pass recovers the competent sequential design when it is
near-optimal and captures the coupling where it is strong (the general
non-separability of Section~\ref{sec:lqg}), without committing to a place-then-control
ordering that can, in more strongly coupled plants, strand the design at a
non-stationary point.

Finally, because the PLL bandwidths and damping gains are \emph{software-settable}
and the useful allocation is operating-point dependent (the localized disturbance
fixes \emph{which} bus is fastest and least certain), the low per-solve cost of
ContEst---one convex solve plus the envelope contraction of
Section~\ref{sec:sim:complexity}---makes it natural to re-run as a slow,
event-triggered design layer that re-provisions set-points as the operating point
shifts, complementary to the online input adaptation of dual/adaptive control.

\begin{table}[t]
\centering
\caption{Low-inertia grid (PLL/DSE): baseline vs.\ co-designed parameters. The
PLL bandwidths $b_i$ enter both $\bar f$ and the measurement covariance
$V(\theta)$.}
\label{tab:pll}
\begin{tabular}{llcc}
\hline
parameter & axis & baseline & optimal \\
\hline
$e_1$ & $f$ (inverter damping) & 1.000 & 3.000 \\
$e_2$ & $f$ (inverter damping) & 1.000 & 1.951 \\
$e_3$ & $f$ (inverter damping) & 1.000 & 2.191 \\
$b_1$ & $h,V$ (PLL bandwidth)  & 6.000 & 13.718 \\
$b_2$ & $h,V$ (PLL bandwidth)  & 6.000 & 10.567 \\
$b_3$ & $h,V$ (PLL bandwidth)  & 6.000 & 10.575 \\
\hline
\end{tabular}
\end{table}

%-----------------------------------------------------------------------
\subsection{Multi-terminal HVDC droop coordination}
\label{sec:sim:mtdc}
%-----------------------------------------------------------------------

The three studies above are nonlinear and use the eKF--MPC inner problem. We close
with a \emph{large-scale, linear} power-system study, and use it to exercise the
\emph{robust output-feedback} ($H_\infty$) inner problem of Section~\ref{sec:hinf} in
place of the $H_2$/LQG one. The device--network model is linearized about an
operating point and discretized by an explicit Euler step,
\begin{equation}\label{eq:lin-ss}
\begin{aligned}
x_{k+1}&=A(\theta)\,x_k+B\,u_k+E\,w_k,\qquad
y_k=C\,x_k+D_v v_k,\\
z_k&=[\,Q^{1/2}x_k;\,R^{1/2}u_k\,],
\end{aligned}
\end{equation}
with $A=I+\Delta t\,A_c$, $B=\Delta t\,B_c$, disturbance channel $E=W^{1/2}$
(renewable infeed), and a measurement $y$ that telemeters only the \emph{DC-bus
voltages} ($C$ selects the $\Delta V_i$ states, $r_y=25$, noise
$D_vD_v^\top=V=10^{-3}I$)---the converter \emph{currents} are unmeasured and must be
estimated. This is therefore a genuine output-feedback co-design. Fixing the
attenuation level $\gamma^2=8$, the inner value splits, as in the LQG separation,
into a worst-case control and a worst-case estimation cost,
\begin{equation}\label{eq:mtdc-Jc}
J_c(\theta)=\underbrace{\operatorname{tr}(XW)}_{J_{\mathrm{det}}(\theta)}
          +\underbrace{\operatorname{tr}(M\Sigma_e)}_{J_{\mathrm{est}}(\theta)},
\end{equation}
with $X$ from the control game Riccati~\eqref{eq:gare} and $\Sigma_e$ from its dual
filter GARE (estimation-error weight $M=Q$); the design gradient is the sum of the
two closed-form contractions~\eqref{eq:gare_grad} and its dual, chained with the
Jacobian of $A(\theta)$---no differentiation through the solver. Although the
fixed-$\theta$ game separates into two Riccatis, the droop $\theta$ enters $A(\theta)$
and hence moves \emph{both} $X$ and $\Sigma_e$, so $J_{\mathrm{det}}$ and
$J_{\mathrm{est}}$ vary over a common $\theta$ and the co-design does not separate.
We take the $O(r_x^3)$ game-Riccati route (Section~\ref{sec:hinf}); the SDP form
\eqref{eq:Hinf_cont}--\eqref{eq:Hinf_est} is avoided because, at the interior
$\gamma$ optimum it minimizes over, its LMIs are rank-deficient and their duals
ill-conditioned (Section~\ref{sec:sim:complexity}). The outer L-BFGS loop of
Section~\ref{sec:algorithm} runs from $5$ starts ($\theta_{\mathrm{nom}}$ plus four
random restarts). The design gradient here sums the control and the \emph{dual
filter} envelope contractions.

\emph{Significance --- why $H_\infty$ fits here.} The failure mode of a droop grid
is a \emph{resonant peak}: the converter lag makes excessive droop provoke a
lightly damped inter-terminal DC-voltage oscillation, and a renewable-infeed step
that lands near that resonance is amplified far beyond its RMS level. It is this
worst-case amplification, not the noise-averaged variance, that a secure design
must bound---exactly the quantity $H_\infty$ prices, while $H_2$ averages it away.
Both halves inherit the resonance: a near-oscillatory plant is both harder to
regulate and harder to \emph{estimate} under a worst-case disturbance, so
$J_{\mathrm{det}}$ and $J_{\mathrm{est}}$ are each sharply $U$-shaped in the droop.
A naive design chasing tight regulation over-provisions droop and lands on that
wall---the baseline here---so ContEst must \emph{lower} the gains to their interior
robust optimum.

\emph{State-space model.} $n_t=25$ terminals on a meshed DC grid ($r_x=50$), each
with state $[\Delta V_i,\Delta I_i]$ and, for $\Delta t=\SI{5}{ms}$,
\begin{subequations}\label{eq:mtdc}
\begin{align}
C_i\,\Delta\dot V_i &= -\!\!\sum_{j\sim i}\!G_{ij}(\Delta V_i-\Delta V_j)+\Delta I_i+\text{infeed},\\
\tau_i\,\Delta\dot I_i &= -\Delta I_i-(k_0\,\theta_i)\,\Delta V_i+u_i,\\
y_{k,i} &= \Delta V_i + v_{k,i},\qquad v_k\sim\mathcal N(0,10^{-3}I_{25}),
\end{align}
\end{subequations}
so the droop gain $k_0\theta_i$ ($k_0=12$) enters $A$ through the
converter-current feedback with lag $\tau_i$. A heterogeneous renewable-infeed
disturbance ($\sigma_{\mathrm{inf}}$ per node) drives the DC nodes; $Q$ penalizes
$\Delta V$, $R=2I$, estimation-error weight $M=Q$; small converter-stress cost
$J_{\mathrm{des}}=c_k\sum_i\theta_i$, $c_k=5\!\times\!10^{-3}$; box
$\theta\in[0.5,4]^{25}$ and naive over-provisioned baseline
$\theta_{\mathrm{nom}}=3$. At this leverage the interior robust optimum is near
$\theta\!\approx\!1$ with $J_c(\theta{=}3)\!\approx\!1.7\,J_c(\theta{=}1)$; only
the DC-bus voltages $\Delta V_i$ are measured, so the converter currents must be
estimated.

\emph{Results.} All five
starts agree, so the optimum is numerically global. Consistent with the $U$-shaped
$J_c$, ContEst \emph{lowers} the over-provisioned droop to an interior,
disturbance-dependent profile ($\theta\in[0.66,0.92]$, none at a bound; noisier
terminals keep slightly higher droop), pulling the terminals off the resonance
wall. Both halves improve: the worst-case estimation cost $J_{\mathrm{est}}$ falls
$42.6\%$ ($1.08\to0.62$) and the worst-case control cost $J_{\mathrm{det}}$ falls
$39.7\%$ ($2.78\to1.68$), for a $40.5\%$ reduction in $J_c$ ($3.86\to2.30$) and,
with the reduced converter stress, a $43.4\%$ reduction in $J_{\mathrm{tot}}$
($4.24\to2.40$). The single control-side design variable improves the estimation
half purely by reshaping the shared dynamics $A(\theta)$---the design-level coupling
that a place-then-control pipeline cannot exploit.

\emph{Adding a sparse sensing axis ($\ell_1$ allocation).} The study above
telemeters \emph{all} $25$ DC-bus voltages. But the meshed ring couples
neighbouring bus voltages, so an un-sensed terminal can be reconstructed by the
$H_\infty$ filter from its neighbours---\emph{which} voltages to telemeter is
itself a design decision, and the natural vehicle for it is the discrete
sensor-allocation relaxation of Section~\ref{sec:sensor}, which the studies so far
have described but not exercised. We give each measurement row a gain
$\alpha_i\ge0$, $C(\alpha)=\operatorname{diag}(\alpha)\,C$ (so $\alpha_i=0$ removes
sensor $i$ and larger $\alpha_i$ raises its SNR at fixed channel noise), append
$\alpha$ to the design, and add the sparsity penalty $\lambda_s\lVert\alpha\rVert_1$
to $J_{\mathrm{des}}$ \eqref{eq:Jdes}. Since $\alpha\ge0$ the epigraph lift of
Section~\ref{sec:sensor} collapses to the smooth linear term
$\lambda_s\mathbf 1^\top\alpha$, and $\alpha$ enters only $C$, whose envelope
gradient \eqref{eq:gradHinf_est} is already in hand---so no machinery changes. This
promotes MTDC to a genuine two-axis co-design: the control droop $\theta_f=k$ (in
$A$) \emph{and} the sensing gains $\theta_h=\alpha$ (in $C$), optimised jointly on
the same robust inner value.

Sweeping $\lambda_s$ traces a sparsity/performance frontier
(Table~\ref{tab:mtdc-alloc}). Across the co-designed rows the control cost
$J_{\mathrm{det}}$ holds at its robust optimum (${\approx}1.67$), so the frontier
isolates the price of sparsity: the worst-case estimation cost degrades
\emph{gracefully} as sensors are pruned, because the network coupling makes most
voltage telemetry redundant---$12$ sensors already leave $J_{\mathrm{est}}$ at the
$25$-sensor \emph{naive}-baseline level, and even $6$ sensors cost only ${\sim}17\%$.
At $\lambda_s=0.06$ the joint co-design keeps $10$ of $25$ sensors; because the
sensor on/off combinatorics make the landscape multimodal (unlike the droop-only
study), we run the multi-start of Section~\ref{sec:algorithm} and report the best of
five. Thresholding that solution to its $10$-sensor integer menu and re-solving the
droop gives $J_{\mathrm{est}}=0.72$ and $J_{\mathrm{det}}=1.67$---$J_c$ down $38.2\%$
against the full-telemetry baseline while \emph{removing $60\%$ of the voltage
sensors}, with estimation in fact \emph{better} than with all $25$ (the pruned
sensors' budget is spent on the survivors). This is the discrete allocation of
Section~\ref{sec:sensor} in action on the robust path: one envelope gradient tunes
the droop and prunes the sensor set at once, a coupling neither a control-only
co-design (no sensing axis) nor a fixed-sensing pipeline (no shared-dynamics
reshaping) can represent.

\begin{table}[t]
\centering
\caption{MTDC sparse voltage-sensor allocation: the $\ell_1$ sparsity/performance
frontier (single-start sweep of $\lambda_s$, joint droop$+$sensing co-design).
Across the co-designed rows $J_{\mathrm{det}}$ holds at its robust optimum, so the
rising $J_{\mathrm{est}}$ isolates the price of removing sensors; the naive
baseline telemeters all $25$ at unit gain with over-provisioned droop $k=3$.}
\label{tab:mtdc-alloc}
\begin{tabular}{lcccc}
\hline
$\lambda_s$ & sensors & $J_{\mathrm{est}}$ & $J_{\mathrm{det}}$ & $J_c$ \\
\hline
baseline & 25 & 1.08 & 2.78 & 3.86 \\
0.00     & 25 & 0.35 & 1.67 & 2.02 \\
0.03     & 15 & 0.80 & 1.67 & 2.47 \\
0.06     & 12 & 1.01 & 1.67 & 2.68 \\
0.10     &  6 & 1.27 & 1.67 & 2.94 \\
0.20     &  2 & 1.53 & 1.66 & 3.20 \\
\hline
\end{tabular}
\end{table}

%-----------------------------------------------------------------------
\subsection{Computational complexity: ContEst with interior-point inner problems}
\label{sec:sim:complexity}
%-----------------------------------------------------------------------

Because the envelope theorem (Section~\ref{sec:envelope}) yields the design
gradient from the \emph{primal and dual of a single inner solve}---never by
differentiating through the solver's iterations---the per-design work is one inner
solve plus a gradient contraction, and the outer L-BFGS iteration inherits the
complexity class of the inner solver. The gradient contraction is cheap: forming
$\partial J/\partial\theta$ from the stored primal/dual costs
$O(r_x^2\,r_\theta)$ (the automatic-differentiation Jacobians of the model
matrices times the dual blocks), dominated by the solve for every case we
consider. The relevant question is therefore the cost of the inner solve.

\emph{Interior-point covariance/LMI solves scale as $O(r_x^6)$.} Both ContEst
inner problems place a covariance matrix among the decision variables. In the LQG
SDPs of Section~\ref{sec:lqg} the variable is the $r_x\times r_x$ covariance
$\Sigma$ (and its dual), i.e.\ $m=O(r_x^2)$ scalar unknowns; a primal--dual
interior-point method assembles and factorizes a dense Schur-complement (Newton)
system of size $m\times m$, at cost $O(m^3)=O(r_x^6)$ per iteration, times the
$O(\sqrt{r_x})$--$O(1)$ iterations interior-point methods take in practice. This is
not the $O(r_x^3)$ of a single matrix factorization: the matrix-valued decision
variable makes the \emph{number} of unknowns quadratic in $r_x$, and it is that
number, cubed, that governs the Newton solve. Empirically, solving the LQR
covariance program with the Clarabel interior-point solver took
$\approx\!2.6\,$s at $r_x=20$ and $\approx\!30\,$s at $r_x=30$---an
$(30/20)^6\!\approx\!11\times$ growth, matching the $O(r_x^6)$ prediction---so a
naive SDP inner problem caps ContEst at a few tens of states.

\emph{The eKF--MPC inner problem meets the same wall.} There the decision trajectory
is the information state $\zeta=[x;\operatorname{vech}\Sigma]$ of dimension
$r_x+\tfrac{r_x(r_x+1)}{2}=O(r_x^2)$; the horizon-$N$ QP \eqref{eq:mpc} thus has
$O(N r_x^2)$ variables and its interior-point solve again grows steeply in $r_x$
(we measured $\approx\!33\,$s per evaluation at $r_x=32$ on the swing model).
Consequently the nonlinear/constrained studies of
Sections~\ref{sec:sim:adcs}--\ref{sec:sim:pll} are run at modest $r_x$, and
large $r_x$ on this path requires model reduction or structure exploitation
(sparsity, chordal decomposition, first-order conic solvers).

\emph{The linear--quadratic case is $O(r_x^3)$ via Riccati.} When the inner problem
is linear and unconstrained---the LQG setting of Section~\ref{sec:lqg}, and, in its
robust form, the MTDC study of Section~\ref{sec:sim:mtdc}---the \emph{same} optimal
value and envelope gradient are available without any interior-point step. The
$H_2$ control cost $J_c=\operatorname{tr}(P W)$ comes from the control Riccati
equation; the gradient needs only the stationary closed-loop covariance $\Sigma_x$
from one discrete Lyapunov solve; and the estimation half comes from the dual
(filter) Riccati. The robust ($H_\infty$) output-feedback value used for MTDC ---
$J_{\mathrm{det}}=\operatorname{tr}(XW)$ and $J_{\mathrm{est}}=\operatorname{tr}(M\Sigma_e)$
--- is identical in cost: both game Riccatis \eqref{eq:gare} are DAREs in augmented
data, so they use the \emph{same} doubling and the \emph{same} envelope identities
\eqref{eq:gare_grad}. Each of these is $O(r_x^3)$, and
structure-preserving doubling converges in a handful of iterations: the Riccati
solve took $\approx\!13\,$ms at $r_x=200$ and $\approx\!60\,$ms at $r_x=400$
(the MTDC game Riccati at $r_x=50$ solves in $\approx\!2\,$ms), i.e.\
three-to-four orders of magnitude faster than the SDP at the crossover and scalable
to the hundreds of states large power-system models require. For the $H_2$ case the
interior-point covariance program and the Riccati recursion return the identical
value and gradient (they agreed to $2\!\times\!10^{-7}$ in our tests); we prefer the
SDP form when its flexibility is needed---indefinite couplings, output-variance or
chance constraints, or the joint estimation--control LMI---and the Riccati form, at
$O(r_x^3)$, whenever the inner problem is a plain (or game-type) linear-quadratic
regulator. The $H_\infty$ case tips the balance further toward Riccati: at the
minimal-$\gamma$ optimum the bounded-real LMI \eqref{eq:Hinf_cont} is rank-deficient
and its dual ill-conditioned, whereas the fixed-$\gamma$ game Riccati is
well-posed---so for robust co-design we use the Riccati form throughout.

\emph{Takeaway.} ContEst's outer cost is (number of designs evaluated) $\times$
(one inner solve $+\;O(r_x^2 r_\theta)$). The envelope gradient keeps the
per-design overhead negligible, so scalability is entirely a property of the inner
solver: $O(r_x^6)$ for interior-point covariance/LMI or information-state MPC
solves (practical to a few tens of states), versus $O(r_x^3)$ for the Riccati
route in the linear-quadratic case (practical to hundreds). Choosing the inner
solver to match the problem structure, rather than defaulting to a general-purpose
conic solver, is what makes large-scale ContEst tractable.

%-----------------------------------------------------------------------
%-----------------------------------------------------------------------
\subsection{Summary of improvements}
\label{sec:sim:summary}
%-----------------------------------------------------------------------

Table~\ref{tab:formulations} consolidates the four studies---their state/input/output
dimensions, the design vector $\theta=(\theta_f,\theta_h)$ and where each axis enters
($\bar f$, $\bar h$, or the noise level $V$), the inner value used, and the
estimation/control/total reductions---and Table~\ref{tab:summary} gives the absolute
cost split for the three nonlinear (eKF--MPC) studies. In those three the total cost
falls by $6$--$35\%$, with a reduction in the estimation
cost $J_{\mathrm{est}}$ in every case, and---by construction---each optimum is an
interior trade rather than a parameter pinned at a bound. The nature of the gain
varies with the physics: for the spacecraft, a single
\emph{power/mass budget} is split across the star tracker, rate gyro and reaction
wheel, favouring attitude sensing ($J_{\mathrm{est}}$ down $23.8\%$, the largest
sensing gain of the closed-loop studies); for the column both axes are
\emph{placements}---feed streams and sensor trays---and the gain is estimation-led
($J_{\mathrm{est}}$ down $22.8\%$ against a nearly flat $J_{\mathrm{det}}$);
and for the low-inertia grid the sensing axis is the \emph{tuning of a dynamic
estimator} (PLL bandwidth), which enters the measurement covariance $V(\theta)$ and,
paired with the extra inverter damping it unlocks, drives by far the largest total
reduction ($34.8\%$). The large-scale MTDC study is reported separately because its
inner value is the robust $H_\infty$ worst-case (guaranteed) cost, on a different
numerical scale from the sampled closed-loop costs of the other three: there a single
\emph{control-side} droop axis reshapes the shared dynamics $A(\theta)$ and thereby
improves \emph{both} halves of the output-feedback design---$J_{\mathrm{est}}$ down
$42.6\%$, $J_{\mathrm{det}}$ down $39.7\%$, $J_{\mathrm{tot}}$ down $43.4\%$
(Table~\ref{tab:formulations})---the design-level coupling a place-then-control
pipeline cannot exploit. Every study compares against the no-upgrade / naive
baseline $\theta_{\mathrm{nom}}$ of its section (for the eKF--MPC studies, a
certainty-equivalence controller on the eKF mean; Section~\ref{sec:numerical}),
at a cost of a handful of inner convex solves per design.

\begin{table*}[t]
\centering
\caption{Consolidated view of the four studies: dimensions $(r_x,r_u,r_y)$; the
design vector $\theta=(\theta_f,\theta_h)$ as $n_f{+}n_h$ with where each axis enters
($\bar f$ dynamics/actuation, $\bar h$ output map, or the measurement covariance
$V$); the inner value (eKF--MPC yields $H_2$/sampled closed-loop costs, MTDC the
$H_\infty$ worst-case guaranteed costs); and the estimation/control/total reductions.
MTDC's single control-side axis still improves $J_{\mathrm{est}}$ because $\theta_f$
reshapes the shared dynamics that the H$_\infty$ filter must estimate through.}
\label{tab:formulations}
\begin{tabular}{lcclcccc}
\hline
study & $(r_x,r_u,r_y)$ & $\theta$: $n_f{+}n_h$ (enters) & inner value
& $J_{\mathrm{est}}\!\downarrow$ & $J_{\mathrm{det}}\!\downarrow$ & $J_{\mathrm{tot}}\!\downarrow$\\
\hline
ADCS (spacecraft) & $(6,3,6)$   & $1{+}2$: $e_{\mathrm{rw}}\!\in\!\bar f$, $\alpha\!\in\!V$ & eKF--MPC ($H_2$) & $23.8\%$ & $4.0\%$  & $\mathbf{7.7\%}$\\
Distillation      & $(12,1,3)$  & $3{+}3$: feeds $\in\!\bar f$, sensors $\in\!\bar h$ & eKF--MPC ($H_2$) & $22.8\%$ & $1.9\%$  & $\mathbf{6.5\%}$\\
PLL/DSE (grid)    & $(9,3,6)$   & $3{+}3$: $e\!\in\!\bar f$, $b\!\in\!\bar f,V$ & eKF--MPC ($H_2$) & $20.9\%$ & $40.6\%$ & $\mathbf{34.8\%}$\\
MTDC (HVDC)       & $(50,25,25)$& $25{+}0$: $k\!\in\!\bar f$ (droop)            & $H_\infty$ output-fb. & $42.6\%$ & $39.7\%$ & $\mathbf{43.4\%}$\\
\hline
\end{tabular}
\end{table*}

\begin{table*}[t]
\centering
\caption{Summary of cost reductions (baseline $\to$ ContEst optimum
$\theta^\star$). $J_{\mathrm{est}}$: estimation cost; $J_{\mathrm{det}}$: control
cost; $J_c=J_{\mathrm{det}}+J_{\mathrm{est}}$;
$J_{\mathrm{tot}}=J_c+J_{\mathrm{des}}$. Percentages are reductions
(larger is better). Baseline is the naive / no-upgrade $\theta_{\mathrm{nom}}$ of
each study (uniform budget for ADCS, naive
default placement for distillation, nominal actuators and PLL for the low-inertia
grid).}
\label{tab:summary}
\resizebox{\textwidth}{!}{%
\begin{tabular}{lccccccccc}
\hline
& \multicolumn{2}{c}{$J_{\mathrm{est}}$}
& \multicolumn{2}{c}{$J_{\mathrm{det}}$}
& \multicolumn{2}{c}{$J_c$}
& $J_{\mathrm{des}}$
& \multicolumn{2}{c}{$J_{\mathrm{tot}}$}\\
\cline{2-3}\cline{4-5}\cline{6-7}\cline{9-10}
study & base $\to$ opt & red. & base $\to$ opt & red. & base $\to$ opt & red.
& opt & base $\to$ opt & red.\\
\hline
ADCS (spacecraft)& $173.1\to132.0$ & $23.8\%$ & $448.3\to430.5$ & $4.0\%$  & $621.4\to562.5$ & $9.5\%$  & $11.0$ & $621.4\to573.5$  & $\mathbf{7.7\%}$\\
Distillation     & $12.8\to9.9$    & $22.8\%$ & $19.5\to19.1$   & $1.9\%$  & $32.3\to29.0$   & $10.2\%$ & $1.2$  & $32.3\to30.2$    & $\mathbf{6.5\%}$\\
PLL/DSE (grid)   & $26.6\to21.0$   & $20.9\%$ & $151.9\to90.2$  & $40.6\%$ & $178.5\to111.2$ & $37.7\%$ & $5.2$  & $178.5\to116.4$  & $\mathbf{34.8\%}$\\
\hline
\end{tabular}%
}
\end{table*}

%=======================================================================
\section{Conclusion}\label{sec:conclusion}
%=======================================================================

We introduced ContEst, a co-design framework that brings state and parameter
estimation into control co-design by optimizing actuation and sensing parameters
against a single stochastic (dual control) objective expressed through the
information state. The framework reduces, in the LQG regime, to a pair of
semidefinite programs and, in the locally linearized regime, to a convex
information-state MPC; sensor allocation enters as a mixed-integer selection with
a sparse $\ell_1$ relaxation. In every case the design gradient is obtained
exactly and cheaply from the inner dual variables by the envelope theorem,
sidestepping implicit differentiation of the solution map and differentiation of
the KKT system. The prevalence of sequential design suggests broad
applicability. Future work includes chance-constrained inner problems,
information states beyond the eKF, and decomposition of the outer problem along
MDO architectures for large interconnected systems.

\begin{ack}
The authors thank colleagues at Argonne National Laboratory for helpful
discussions.
\end{ack}

\appendix

%=======================================================================
\section{Proof of the LQG regularity Proposition}\label{app:lqgreg}
%=======================================================================

This appendix proves Proposition~\ref{prop:lqgreg} and derives the dual
identities $S_{11}^\star=P$, $S_{11}^{\prime\star}=\Sigma^\varepsilon$, replacing
the pointer to the sensitivity literature in Corollary~\ref{cor:lqg} with a
self-contained argument. We give the control SDP~\eqref{eq:Jcont}; the estimation
SDP~\eqref{eq:Jest} is its formal dual (swap $A_\theta\!\leftrightarrow
A_\theta^\top$, $B_\theta\!\leftrightarrow C_\theta^\top$, $W\!\leftrightarrow Q$,
$R\!\leftrightarrow V$) and the same steps apply verbatim, yielding
$P^{\varepsilon\star}=\Sigma^\varepsilon$ and $S_{11}^{\prime\star}=
\Sigma^\varepsilon$.

\emph{(a) Slater.} Pick any stabilizing $K_0$ (one exists by stabilizability) and
let $\Sigma_0\succ0$ solve the strict Lyapunov inequality $A_{\mathrm{cl},0}
\Sigma_0 A_{\mathrm{cl},0}^\top-\Sigma_0+W\prec0$ with $A_{\mathrm{cl},0}=
A_\theta+B_\theta K_0$; such $\Sigma_0$ exists because $A_{\mathrm{cl},0}$ is
Schur and $W\succ0$. Setting $L_0=K_0\Sigma_0$ and $Z_0\succ K_0\Sigma_0K_0^\top$
makes both LMIs of \eqref{eq:Jcont} strictly feasible, so Slater holds and strong
duality with attainment follows \cite{boyd2004convex}.

\emph{(b) Primal uniqueness.} With $W\succ0$, the second LMI of \eqref{eq:Jcont}
forces $\Sigma\succ0$ and, taking Schur complements, encodes
$A_{\mathrm{cl}}\Sigma A_{\mathrm{cl}}^\top-\Sigma+W\preceq0$ for the gain
$K=L\Sigma^{-1}$; a Lyapunov-eigenvector argument then shows any feasible
$(\Sigma,L)$ has Schur $A_{\mathrm{cl}}$, i.e.\ encodes a stabilizing gain. The
objective is the stationary cost $J(K)=\operatorname{tr}(Q\Sigma_K)+
\operatorname{tr}(RK\Sigma_KK^\top)$, and completion of squares against the
optimal $K^\star$ gives
\begin{equation}\label{eq:cos}
J(K)-\operatorname{tr}(PW)=\operatorname{tr}\!\big((K-K^\star)^\top
(R+B_\theta^\top PB_\theta)(K-K^\star)\Sigma_K\big),
\end{equation}
with $P$ the stabilizing solution of the control DARE. Since $R+B_\theta^\top
PB_\theta\succ0$ and $\Sigma_K\succ0$, the right-hand side is zero iff
$K=K^\star$, so the minimizer $\Sigma^\star=\Sigma_{K^\star}$ (and
$L^\star=K^\star\Sigma^\star$) is unique.

\emph{(c) Dual uniqueness, strict complementarity, and $S_{11}^\star=P$.} Let
$S=\left[\begin{smallmatrix}S_{11}&S_{12}\\S_{12}^\top&S_{22}\end{smallmatrix}
\right]\succeq0$ be the dual of the second LMI. Complementary slackness aligns the
range of $S$ with the kernel of the (rank-deficient at optimum) primal slack,
$\ker\mathcal G^\star=\operatorname{range}[\,I;\,-A_{\mathrm{cl}}^{\star\top}\,]$,
which parametrizes $S$ through $S_{12}=-S_{11}A_{\mathrm{cl}}^\star$ (so
$S_{12}^\star=-PA_{\mathrm{cl}}^\star$) up to the free block
$S_{11}$. Stationarity of the Lagrangian in $\Sigma$ then reads
\begin{equation}\label{eq:stein}
S_{11}=Q+K^{\star\top}RK^\star+A_{\mathrm{cl}}^\top S_{11}A_{\mathrm{cl}},
\end{equation}
the closed-loop Stein (Lyapunov) equation, whose solution is unique because
$A_{\mathrm{cl}}$ is Schur. But the control DARE is exactly $P=Q+
K^{\star\top}RK^\star+A_{\mathrm{cl}}^\top PA_{\mathrm{cl}}$, so \eqref{eq:stein}
gives $S_{11}=P$: the $(1,1)$ dual block \emph{is} the control Riccati solution,
and the dual is unique. Strict complementarity holds because $P\succeq Q\succ0$
and $R\succ0$ make the primal and dual slacks have complementary full rank. This
verifies Assumption~\ref{as:env}(A1)--(A3); (A4) is the modeling hypothesis that
$\theta$ enters only through $C^1$ data, and Corollary~\ref{cor:lqg} follows.
\hfill$\square$

The identity $S_{11}^\star=P$ (and dually $S_{11}^{\prime\star}=
\Sigma^\varepsilon$) is what makes the covariance-SDP gradient
$\partial J_{\mathrm{cont}}/\partial W=-S_{11}^\star$ of \eqref{eq:gradQRcont}
numerically identical to the Riccati gradient $\partial\operatorname{tr}(PW)/
\partial W=P$ used in Section~\ref{sec:sim:complexity}.

%=======================================================================
\section{Proof of the fixed-$\gamma$ $H_\infty$ regularity Proposition}
\label{app:hinfreg}
%=======================================================================

This appendix proves Proposition~\ref{prop:hinfreg}, the robust twin of
Appendix~\ref{app:lqgreg}: at a fixed attenuation $\gamma$ above the optimal
one, the game GARE~\eqref{eq:gare} plays exactly the role the DARE plays in the
LQG case, so its stabilizing solution is a locally unique $C^1$ function of
$\theta$ and $V_{\mathrm{wc}}(\theta)=\operatorname{tr}(XW_\theta)$ carries the
closed-form gradient~\eqref{eq:gare_grad}. The argument parallels
Appendix~\ref{app:lqgreg} step for step, the only change being that the
completion-of-squares now carries the \emph{indefinite} game weight
$\tilde R=\operatorname{diag}(R,-\gamma^2 I)$; positive-definiteness of $R$ is
nowhere used, exactly as anticipated in Section~\ref{sec:hinf}. We treat the
control GARE; the filter GARE follows by the transpose duality of
Appendix~\ref{app:lqgreg}. Recall $\tilde B_\theta=[\,B_\theta\ \ E_\theta\,]$,
$E_\theta=W_\theta^{1/2}$, the saddle gain
$\tilde K=-(\tilde R+\tilde B_\theta^\top X\tilde B_\theta)^{-1}
\tilde B_\theta^\top X A_\theta$, and $A_{\mathrm{cl}}=A_\theta+\tilde B_\theta\tilde K$.

\emph{(a) Stabilizing solution.} For $\gamma$ above the optimal attenuation
$\gamma^\star$, the discrete dynamic-game (equivalently, discrete bounded-real)
Riccati theory gives a unique $X\succeq0$ solving \eqref{eq:gare} for which
$R+B_\theta^\top XB_\theta\succ0$, the full game pivot $\tilde R+\tilde
B_\theta^\top X\tilde B_\theta$ is nonsingular with inertia $(r_u,r_w)$, and the
game closed loop $A_{\mathrm{cl}}$ is Schur \cite{basar1995dynamic,boyd1994linear}.
This is the fixed-$\gamma$ analog of the classical DARE existence invoked in
Appendix~\ref{app:lqgreg}, and it is exactly the admissibility used in
Section~\ref{sec:hinf}. (As $\gamma^\star$ is approached the pivot loses rank,
which is the rank drop flagged in Remark~\ref{rem:hinfreg}; we stay strictly above
it.)

\emph{(b) Lyapunov (completion-of-squares) form.} Substituting
$\tilde K$ into $Q+A_{\mathrm{cl}}^\top XA_{\mathrm{cl}}+\tilde K^\top\tilde
R\tilde K$ and cancelling gives
\begin{equation}\label{eq:gare-lyap}
X=Q+\tilde K^\top\tilde R\tilde K+A_{\mathrm{cl}}^\top XA_{\mathrm{cl}},
\end{equation}
which is identical algebra to the LQG identity below \eqref{eq:stein}, valid for
the indefinite $\tilde R$ because completion of squares never uses its sign. This
is the robust twin of $P=Q+K^{\star\top}RK^\star+A_{\mathrm{cl}}^\top
PA_{\mathrm{cl}}$.

\emph{(c) Local uniqueness and $C^1$ dependence.} Write the GARE residual
$\mathcal R(X,\theta)=Q+A_\theta^\top XA_\theta-A_\theta^\top X\tilde B_\theta
(\tilde R+\tilde B_\theta^\top X\tilde B_\theta)^{-1}\tilde B_\theta^\top
XA_\theta-X$. It is rational and $C^1$ in $(X,\theta)$ wherever the pivot is
nonsingular, hence near a stabilizing root by (a). Because $\tilde K$ is the
saddle point of the inner game (nonsingular pivot), the envelope of the game
kills the multiplier's first-order variation and the partial Fr\'echet derivative
in $X$ collapses to the closed-loop Stein operator,
\begin{equation}\label{eq:gare-ift}
\partial_X\mathcal R(X,\theta)[\delta X]
=A_{\mathrm{cl}}^\top\,\delta X\,A_{\mathrm{cl}}-\delta X,
\end{equation}
which is invertible precisely because $A_{\mathrm{cl}}$ is Schur (its eigenvalues
are $\lambda_i\bar\lambda_j-1\neq0$ for $|\lambda_i|<1$). The implicit function
theorem then yields a unique $C^1$ branch $\theta\mapsto X(\theta)$ through the
stabilizing root, and with it $C^1$ maps $\tilde K(\theta)$, $A_{\mathrm{cl}}(\theta)$,
and the game closed-loop gramian $S(\theta)$ solving
$S=A_{\mathrm{cl}}SA_{\mathrm{cl}}^\top+W_\theta$. This is the content of
Proposition~\ref{prop:hinfreg}; it is the direct analog of
Lemma-level smoothness of the DARE solution in Appendix~\ref{app:lqgreg}, and it
needs neither an SDP dual nor strict complementarity.

\emph{(d) $C^1$ value and the closed-form gradient.} Since
$V_{\mathrm{wc}}=\operatorname{tr}(XW_\theta)$ with $X\in C^1$ and $W_\theta\in
C^1$, $V_{\mathrm{wc}}\in C^1$. Differentiating \eqref{eq:gare-lyap}, the saddle
stationarity of $\tilde K$ annihilates all $\mathrm d\tilde K$ terms, so
$\mathrm dX$ solves the Stein equation $\mathrm
dX-A_{\mathrm{cl}}^\top\mathrm dX A_{\mathrm{cl}}=\Xi$ with source
$\Xi=\mathrm dQ+\tilde K^\top\mathrm d\tilde R\,\tilde K+
A_{\mathrm{cl}}^\top X(\mathrm dA_\theta+\mathrm d\tilde B_\theta\,\tilde K)
+(\cdot)^\top$. Contracting with the gramian $S$ and using
$\operatorname{tr}(W_\theta\,\mathrm dX)=\operatorname{tr}\!\big(S(\mathrm
dX-A_{\mathrm{cl}}^\top\mathrm dX A_{\mathrm{cl}})\big)=\operatorname{tr}(S\Xi)$
gives, term by term,
\begin{equation}\label{eq:gare-grad-app}
\begin{aligned}
&\frac{\partial V_{\mathrm{wc}}}{\partial A_\theta}=2\,XA_{\mathrm{cl}}S,\qquad
\frac{\partial V_{\mathrm{wc}}}{\partial \tilde B_\theta}=2\,XA_{\mathrm{cl}}S\tilde K^\top,\\
&\frac{\partial V_{\mathrm{wc}}}{\partial Q}=S,\qquad
\frac{\partial V_{\mathrm{wc}}}{\partial \tilde R}=\tilde KS\tilde K^\top,
\end{aligned}
\end{equation}
together with the explicit $\partial V_{\mathrm{wc}}/\partial W_\theta=X$ from the
objective; these are exactly \eqref{eq:gare_grad}. The disturbance channel
$E_\theta=W_\theta^{1/2}$ enters through the $E$-block of
$\partial V_{\mathrm{wc}}/\partial\tilde B_\theta$, so when a single design
coordinate $\theta_i$ moves both the objective weight $W_\theta$ and the channel
$E_\theta$ its gradient sums the two contributions, precisely the data-object
chain rule \eqref{eq:chain} of Assumption~\ref{as:env}(A4).

\emph{(e) Filter GARE by duality.} Under the transpose substitution of
Appendix~\ref{app:lqgreg}---$A_\theta\!\mapsto\!A_\theta^\top$, $\tilde
B_\theta\!\mapsto\![\,C_\theta^\top\ L_e^\top\,]$, $\tilde
R\!\mapsto\!\operatorname{diag}(V,-\gamma^2 I)$, $W_\theta$ the driving
covariance---the worst-case a-priori error covariance $\Sigma_e$ solves the
control GARE on the dual data, so (a)--(d) transfer verbatim: $\Sigma_e(\theta)$
is a locally unique $C^1$ stabilizing solution and
$V_{\mathrm{est}}(\theta)=\operatorname{tr}(M\Sigma_e)$ has the dual gradient
(contract $2\,\Sigma_e\tilde A_{\mathrm{cl}}\tilde S$ against
$\partial A_\theta^\top/\partial\theta$, with the $C_\theta,V$ blocks of
$2\,\Sigma_e\tilde A_{\mathrm{cl}}\tilde S\tilde K^\top$), matching
Section~\ref{sec:hinf}.

\emph{(f) The $H_2$ limit.} As $\gamma\!\to\!\infty$ the disturbance player is
switched off: the $-\gamma^2 I$ pivot block deactivates, the GARE~\eqref{eq:gare}
reduces to the control DARE $P=Q+K^{\star\top}RK^\star+A_{\mathrm{cl}}^\top
PA_{\mathrm{cl}}$ of Appendix~\ref{app:lqgreg}, $X\!\to\!P$,
$S\!\to\!\Sigma^\star$, and \eqref{eq:gare-grad-app} collapses to the LQG
gradients of Appendix~\ref{app:lqgreg}. The $H_2$/LQG result is thus the
$\gamma\!\to\!\infty$ boundary of this one, and (unlike the minimal-$\gamma$ BRL
SDP of Remark~\ref{rem:hinfreg}) no rank drop or ill-conditioned dual ever enters
the fixed-$\gamma$ route. \hfill$\square$

\bibliographystyle{plain}
\bibliography{References}

% ----------------------------------------------------------------------
% Suggested BibTeX for two works you already own (verify before use):
%
% @inproceedings{grigoriadis1998integrate,
%   author    = {Grigoriadis, Karolos M. and Zhu, Guoming and Skelton, Robert E.},
%   title     = {Optimal Redesign of Linear Systems / An Algorithm to Integrate
%                Structure and Control Design},
%   booktitle = {Proc.\ IEEE Conf.\ Decision and Control (CDC)},
%   year      = {1996}}   % VERIFY title, venue, year, pages
%
% @article{ramadan_ekfkoopman,
%   author  = {Ramadan, Mohammad S. and Anitescu, Mihai},
%   title   = {An eKF--Koopman Operator Approach for Tractable Stochastic
%              Optimal Control},
%   journal = {IEEE ...},
%   year    = {2024}}      % VERIFY title, journal, volume, pages, year
% ----------------------------------------------------------------------

\end{document}